\NormalFont\ShortTitle{Faptele}

{\MT Faptele Apostolilor


\par }\ChapOne{1}{\PP \VerseOne{1}Prima carte pe care am scris-o, Teofilule, se referă la tot ce a început Isus să facă și să învețe,

\VS{2}până în ziua în care a fost înălțat, după ce a dat poruncă prin Duhul Sfânt apostolilor pe care îi alesese.

\VS{3}Acestora li s-a arătat viu și după ce a suferit, prin multe dovezi, arătându-li-se pe parcursul a patruzeci de zile și vorbindu-le despre Împărăția lui Dumnezeu.

\VS{4}Fiind adunat împreună cu ei, le-a poruncit:
{\WJ{“Nu plecați din Ierusalim, ci așteptați promisiunea Tatălui, pe care ați auzit-o de la mine.
}}

\VS{5}{\WJ{Căci Ioan a botezat într-adevăr în apă, dar voi veți fi botezați în Duhul Sfânt nu peste multe zile”.
}}


\par }{\PP \VS{6}Când s-au adunat, L-au întrebat: “Doamne, acum dai înapoi împărăția lui Israel?”


\par }{\PP \VS{7}El le-a zis: “{\WJ{Nu este treaba voastră să cunoașteți vremurile și timpurile pe care Tatăl le-a pus în autoritatea Sa.
}}

\VS{8}{\WJ{Dar voi veți primi putere când va veni Duhul Sfânt peste voi. Îmi veți fi martori în Ierusalim, în toată Iudeea și Samaria și până la marginile pământului.”
}}


\par }{\PP \VS{9}După ce a spus aceste lucruri, pe când se uitau ei, a fost ridicat și un nor L-a luat din fața lor.

\VS{10}În timp ce ei priveau neclintit spre cer, în timp ce El se ducea, iată că lângă ei stăteau doi bărbați îmbrăcați în haine albe,

\VS{11}care au spus și ei: “Bărbați galileeni, de ce stați cu ochii în cer? Acest Isus, care a fost primit de la voi în cer, se va întoarce în același mod în care L-ați văzut mergând în cer.”


\par }{\PP \VS{12}Apoi s-au întors la Ierusalim, de pe muntele numit Muntele Măslinilor, care este lângă Ierusalim, la o zi de drum de Sabat.

\VS{13}După ce au ajuns, s-au urcat în camera de sus unde stăteau: Petru, Ioan, Iacov, Iacov, Andrei, Filip, Toma, Bartolomeu, Matei, Iacov, fiul lui Alfeu, Simon Zelota și Iuda, fiul lui Iacov.

\VS{14}Toți aceștia, cu un singur gând, stăruiau cu stăruință în rugăciune și în cereri, împreună cu femeile, cu Maria, mama lui Isus, și cu frații Lui.


\par }{\PP \VS{15}În zilele acelea, Petru s-a ridicat în picioare în mijlocul ucenicilor, care erau în număr de o sută douăzeci, și a zis:

\VS{16}“Fraților, trebuia să se împlinească Scriptura aceasta, pe care Duhul Sfânt a spus-o mai înainte, prin gura lui David, despre Iuda, care era călăuza celor ce luaseră pe Isus.

\VS{17}Căci el a fost socotit împreună cu noi și a primit partea lui în această slujbă.

\VS{18}Omul acesta a obținut un câmp cu răsplata pentru răutatea lui; și, căzând cu capul în jos, i s-a deschis trupul și i-au țâșnit toate intestinele.

\VS{19}Toți cei care locuiau în Ierusalim au aflat că în limba lor câmpul acela se numea “Akeldama”, adică “Câmpul de sânge”.

\VS{20}Căci este scris în cartea Psalmilor: “Este un câmp de sânge,

\par }{\Q “Să fie pustiită locuința lui.

\par }{\QB Nimeni să nu locuiască în ea”.

\par }{\PP și,

\par }{\Q “Lasă-l pe altul să își ocupe funcția.


\par }{\PP \VS{21}“Așadar, dintre bărbații care ne-au însoțit în tot timpul în care Domnul Isus a intrat și a ieșit dintre noi,

\VS{22}începând de la botezul lui Ioan și până în ziua în care a fost ridicat de la noi, unul dintre aceștia trebuie să fie martor cu noi al învierii Lui.”


\par }{\PP \VS{23}Au propus două: Iosif, zis Barsabas, care se mai numea și Iustus, și Matia.

\VS{24}Ei s-au rugat și au zis: “Tu, Doamne, care cunoști inimile tuturor oamenilor, arată pe care dintre acești doi l-ai ales

\VS{25}pentru a lua parte la această slujbă și apostolat, de la care a căzut Iuda, ca să se ducă la locul lui.”

\VS{26}Au tras la sorți pentru ei și sorțul a căzut pe Matia; și a fost numărat împreună cu cei unsprezece apostoli.



\par }\Chap{2}{\PP \VerseOne{1}Când a sosit ziua Cincizecimii, erau toți la un loc, deodată, în același loc.

\VS{2}Deodată s-a auzit din cer un zgomot ca un vuiet de vânt puternic, care a umplut toată casa în care stăteau.

\VS{3}Au apărut limbi ca de foc și li s-au împărțit, și câte una s-a așezat pe fiecare dintre ei.

\VS{4}Cu toții au fost umpluți de Duhul Sfânt și au început să vorbească în alte limbi, după cum le dădea Duhul să vorbească.


\par }{\PP \VS{5}În Ierusalim locuiau iudei evrei, oameni evlavioși, din toate neamurile de sub cer.

\VS{6}Când s-a auzit acest sunet, mulțimea s-a adunat și a rămas nedumerită, pentru că fiecare îi auzea vorbind în limba lui.

\VS{7}Toți erau uimiți și se mirau, spunându-și unii altora: “Iată, oare nu sunt galileeni toți aceștia care vorbesc?

\VS{8}Cum îi auzim, fiecare în limba sa maternă?

\VS{9}Parțieni, medi, elamiți și oameni din Mesopotamia, Iudeea, Capadocia, Pont, Asia,

\VS{10}Frigia, Pamfilia, Egipt, părțile Libiei din jurul Cirenei, vizitatori din Roma, atât evrei cât și prozeliți,

\VS{11}cretani și arabi — noi îi auzim vorbind în limbile noastre despre lucrările mărețe ale lui Dumnezeu!”

\VS{12}Toți au rămas uimiți și perplecși, spunându-și unii altora: “Ce înseamnă aceasta?”

\VS{13}Alții, batjocorind, spuneau: “Sunt plini de vin nou”.


\par }{\PP \VS{14}Dar Petru, stând în picioare cu cei unsprezece, a ridicat glasul și le-a zis: “Bărbați din Iudeea și voi toți cei ce locuiți la Ierusalim, să vă fie cunoscut lucrul acesta și ascultați cuvintele mele.

\VS{15}Căci aceștia nu sunt beți, așa cum credeți voi, întrucât este abia al treilea ceas al zilei.

\VS{16}Ci iată ce s-a spus prin profetul Ioel:


\par }{\Q \VS{17}“Așa va fi în zilele de pe urmă, zice Dumnezeu,

\par }{\QB că Îmi voi revărsa Duhul Meu peste toată făptura.

\par }{\Q Fiii și fiicele voastre vor profeți.

\par }{\QB Tinerii tăi vor avea viziuni.

\par }{\QB Bătrânii tăi vor visa vise.


\par }{\Q \VS{18}Și pe robii mei și pe roabele mele în zilele acelea,

\par }{\QB Voi revărsa Duhul Meu și ei vor profeți.


\par }{\Q \VS{19}Voi face minuni pe cerul de sus,

\par }{\QB și semne pe pământul de dedesubt:

\par }{\QB sânge, foc și valuri de fum.


\par }{\Q \VS{20}Soarele se va preface în întuneric,

\par }{\QB și luna în sânge,

\par }{\QB înainte ca ziua cea mare și glorioasă a Domnului să vină.


\par }{\Q \VS{21}Oricine va invoca Numele Domnului va fi mântuit.


\par }{\PP \VS{22}“Bărbați ai lui Israel, ascultați aceste cuvinte! Pe Isus din Nazaret, un om aprobat de Dumnezeu pentru voi prin fapte puternice, minuni și semne pe care Dumnezeu le-a făcut prin el printre voi, după cum știți voi înșivă,

\VS{23}pe El, care a fost predat prin sfatul hotărât și prin știința mai dinainte a lui Dumnezeu, L-ați prins prin mâna unor oameni fără de lege, L-ați răstignit și L-ați ucis;

\VS{24}pe care Dumnezeu L-a înviat, eliberându-L de agonia morții, pentru că nu era cu putință ca El să fie ținut de ea.

\VS{25}Căci David spune despre El: “Nu este posibil să fi fost învins de Dumnezeu, pentru că nu era cu putință să fie învins,

\par }{\Q 'Îl vedeam pe Domnul mereu înaintea feței mele,

\par }{\QB Căci El este la dreapta mea, ca să nu mă clatin.


\par }{\Q \VS{26}De aceea inima mea s-a bucurat și limba mea s-a veselit.

\par }{\QB Mai mult, și carnea mea va locui în speranță,


\par }{\Q \VS{27}pentru că nu vei lăsa sufletul meu în Locuința morților,

\par }{\QB nici nu vei permite ca Sfântul tău să vadă putreziciunea.


\par }{\Q \VS{28}Tu mi-ai făcut cunoscute căile vieții.

\par }{\QB Mă vei umple de bucurie cu prezența ta.


\par }{\PP \VS{29}“Fraților, pot să vă spun, în mod liber, despre patriarhul David, că a murit și a fost îngropat, iar mormântul lui este la noi până în ziua de azi.

\VS{30}De aceea, fiind profet și știind că Dumnezeu îi jurase cu jurământ că din rodul trupului său, după trup, va ridica pe Hristos ca să șadă pe tronul său,

\VS{31}el, prevăzând acest lucru, a vorbit despre învierea lui Hristos, că sufletul lui nu a rămas în Hades și că trupul lui nu a văzut putrezirea.

\VS{32}Acest Isus pe care Dumnezeu l-a înviat, despre care noi toți suntem martori.

\VS{33}De aceea, fiind înălțat la dreapta lui Dumnezeu și primind de la Tatăl făgăduința Duhului Sfânt, a revărsat ceea ce vedeți și auziți acum.

\VS{34}Căci David nu s-a înălțat la ceruri, ci el însuși spune,

\par }{\Q 'Domnul a zis Domnului meu: “Șezi la dreapta Mea.


\par }{\QB \VS{35}până când voi face din vrăjmașii tăi un tăpșan pentru picioarele tale.”


\par }{\PP \VS{36}“Să știe deci toată casa lui Israel că Dumnezeu a făcut Domn și Hristos pe acest Isus pe care voi L-ați răstignit.”


\par }{\PP \VS{37}Auzind acestea, s-au cutremurat și au zis lui Petru și celorlalți apostoli: “Fraților, ce să facem?”


\par }{\PP \VS{38}Petru le-a zis: “Pocăiți-vă și fiecare dintre voi să fie botezat în numele lui Isus Hristos, pentru iertarea păcatelor, și veți primi darul Duhului Sfânt.

\VS{39}Căci făgăduința este pentru voi și pentru copiii voștri și pentru toți cei de departe, pentru toți cei pe care îi va chema la sine Domnul Dumnezeul nostru.”

\VS{40}Cu multe alte cuvinte a dat mărturie și i-a îndemnat, zicând: “Salvați-vă de acest neam strâmb!”.


\par }{\PP \VS{41}Și cei ce au primit cu bucurie cuvântul Lui au fost botezați. În ziua aceea s-au adăugat în jur de trei mii de suflete.

\VS{42}Ei stăruiau cu stăruință în învățătura apostolilor și în părtășia cu ei, în frângerea pâinii și în rugăciune.

\VS{43}Înfricoșarea a cuprins fiecare suflet, și multe minuni și semne se făceau prin apostoli.

\VS{44}Toți cei care credeau erau împreună și aveau toate lucrurile în comun.

\VS{45}Își vindeau averile și bunurile și le împărțeau la toți, după cum avea cineva nevoie.

\VS{46}Zi de zi, stăruind cu stăruință și de comun acord în templu și frângând pâinea acasă, își luau hrana cu bucurie și cu inimă curată,

\VS{47}lăudând pe Dumnezeu și având parte de bunăvoința întregului popor. Domnul adăuga zi de zi la adunare pe cei care erau salvați.



\par }\Chap{3}{\PP \VerseOne{1}Petru și Ioan se urcau în templu la ora de rugăciune, la ceasul al nouălea.

\VS{2}Se ducea un om șchiop din pântecele mamei sale, pe care îl puneau zilnic la ușa templului numit Frumos, ca să ceară daruri pentru cei nevoiași dintre cei care intrau în templu.

\VS{3}Văzându-i pe Petru și pe Ioan pe cale să intre în templu, a cerut să primească daruri pentru cei nevoiași.

\VS{4}Petru, fixându-și ochii asupra lui, împreună cu Ioan, a zis: “Uită-te la noi”.

\VS{5}El i-a ascultat, așteptând să primească ceva de la ei.

\VS{6}Dar Petru a spus: “Nu am nici argint, nici aur, ci ceea ce am, asta vă dau. În numele lui Isus Hristos din Nazaret, ridică-te și umblă!”.

\VS{7}El l-a luat de mâna dreaptă și l-a ridicat. Imediat picioarele și oasele gleznelor lui au primit putere.

\VS{8}Sărind, s-a ridicat în picioare și a început să meargă. A intrat cu ei în templu, mergând, sărind și lăudându-l pe Dumnezeu.

\VS{9}Tot poporul l-a văzut mergând și lăudându-l pe Dumnezeu.

\VS{10}L-au recunoscut, că era cel care obișnuia să stea să cerșească daruri pentru nevoiași la Poarta Frumoasă a templului. Erau plini de uimire și de mirare pentru ceea ce i se întâmplase.

\VS{11}În timp ce șchiopul vindecat se ținea de Petru și de Ioan, tot poporul a alergat împreună la ei în pridvorul care se numește al lui Solomon, minunându-se foarte mult.


\par }{\PP \VS{12}Petru, văzând aceasta, a zis norodului: “Bărbați israeliți, de ce vă minunați de omul acesta? De ce vă holbați ochii la noi, ca și cum prin puterea sau evlavia noastră l-am fi făcut să meargă?

\VS{13}Dumnezeul lui Avraam, al lui Isaac și al lui Iacov, Dumnezeul părinților noștri, a proslăvit pe Robul Său Isus, pe care voi L-ați predat și L-ați renegat în fața lui Pilat, când acesta hotărâse să-L elibereze.

\VS{14}Dar voi L-ați renegat pe Cel Sfânt și Drept și ați cerut să vi se acorde un ucigaș,

\VS{15}și ați ucis pe Prințul vieții, pe care Dumnezeu L-a înviat din morți, lucru despre care noi suntem martori.

\VS{16}Prin credința în numele lui, numele lui l-a făcut puternic pe acest om, pe care îl vedeți și îl cunoașteți. Da, credința care este prin el i-a dat această desăvârșită sănătate în prezența voastră a tuturor.


\par }{\PP \VS{17}“Fraților,\FTNT{*}{{\FR 3:17 }{\FT{Cuvântul “frați” poate fi tradus corect și prin “frați și surori” sau “frați și surori”.}}} știu că voi, ca și căpeteniile voastre, ați făcut aceasta din neștiință.

\VS{18}Dar lucrurile pe care le-a anunțat Dumnezeu prin gura tuturor profeților Săi, că Hristos va suferi, le-a împlinit astfel.


\par }{\PP \VS{19}“Pocăiți-vă, deci, și întoarceți-vă, ca să vi se șteargă păcatele, ca să vină vremuri de răcorire de la fața Domnului,

\VS{20}ca să trimită pe Hristos Isus, care a fost rânduit pentru voi mai înainte,

\VS{21}pe care cerul trebuie să-L primească până la vremurile de refacere a tuturor lucrurilor, pe care Dumnezeu le-a spus demult prin gura sfinților Săi prooroci.

\VS{22}Căci Moise a spus într-adevăr părinților: “Domnul Dumnezeu vă va ridica un profet dintre frații voștri, ca mine. Îl veți asculta în toate lucrurile pe care vi le va spune.

\VS{23}Orice suflet care nu va asculta de acel profet va fi nimicit cu desăvârșire din mijlocul poporului'.
\FTNT{†}{{\FR 3:23 }{\FT{Deuteronom 18:15,18-19
}}}

\VS{24}Da, și toți profeții, de la Samuel și cei care au urmat după el, toți cei care au vorbit, au vorbit și ei despre aceste zile.

\VS{25}Voi sunteți copiii profeților și ai legământului pe care l-a făcut Dumnezeu cu părinții noștri, când a zis lui Avraam: 'Toate familiile pământului vor fi binecuvântate prin urmașii tăi'.
\FTNT{‡}{{\FR 3:25 }{\FT{sau, semințe}}}\FTNT{§}{{\FR 3:25 }{\FT{Geneza 22:18; 26:4}}}

\VS{26}Dumnezeu, după ce l-a ridicat pe robul său Isus, l-a trimis mai întâi la voi ca să vă binecuvânteze, întorcându-vă pe fiecare dintre voi de la răutatea voastră.”



\par }\Chap{4}{\PP \VerseOne{1}Pe când vorbeau ei poporului, au venit la ei preoții, căpetenia templului și saducheii,

\VS{2}care erau supărați că ei învățau poporul și propovăduiau în Iisus învierea din morți.

\VS{3}Au pus mâna pe ei și i-au pus în arest până a doua zi, căci era deja seară.

\VS{4}Dar mulți dintre cei care au auzit cuvântul au crezut și numărul lor a ajuns să fie cam cinci mii.


\par }{\PP \VS{5}Dimineața, s-au adunat la Ierusalim conducătorii, bătrânii și cărturarii lor.

\VS{6}Marele preot Ana era acolo, împreună cu Caiafa, cu Ioan, cu Alexandru și cu toți cei care erau rude cu marele preot.

\VS{7}După ce au așezat pe Petru și pe Ioan în mijlocul lor, i-au întrebat: “Cu ce putere sau în ce nume ați făcut aceasta?”


\par }{\PP \VS{8}Atunci Petru, plin de Duhul Sfânt, le-a zis: “Căpeteniile poporului și bătrânii lui Israel,

\VS{9}dacă suntem cercetați astăzi pentru o faptă bună făcută unui om olog, și pentru că omul acesta a fost vindecat,

\VS{10}să știți voi toți și tot poporul lui Israel că, în numele lui Isus Hristos din Nazaret, pe care L-ați răstignit și pe care Dumnezeu L-a înviat din morți, omul acesta stă aici înaintea voastră, întreg în el.

\VS{11}El este “piatra care a fost considerată de voi, zidarii, ca fiind fără valoare, și care a devenit capul unghiului”.
\FTNT{*}{{\FR 4:11 }{\FT{{\CHAR{Psalmul 118:22}}}}}

\VS{12}În nimeni altcineva nu este mântuire, pentru că nu există sub cer un alt nume sub cer, dat între oameni, prin care să fim mântuiți!”


\par }{\PP \VS{13}Când au văzut îndrăzneala lui Petru și a lui Ioan, și au înțeles că erau oameni neînvățați și neștiutori, s-au mirat. Ei au recunoscut că fuseseră cu Isus.

\VS{14}Văzând că omul care fusese vindecat stătea cu ei, nu au putut spune nimic împotrivă.

\VS{15}Dar, după ce le-au poruncit să iasă deoparte din consiliu, s-au sfătuit între ei,

\VS{16}spunând: “Ce să le facem acestor oameni? Pentru că, într-adevăr, s-a făcut prin ei o minune remarcabilă, după cum pot vedea cu ochiul liber toți cei care locuiesc în Ierusalim, și nu putem nega acest lucru.

\VS{17}Dar, pentru ca acest lucru să nu se răspândească mai departe în popor, să-i amenințăm, ca de acum înainte să nu mai vorbească nimănui în acest nume.”

\VS{18}Ei i-au chemat și le-au poruncit să nu vorbească deloc și să nu învețe în numele lui Isus.


\par }{\PP \VS{19}Dar Petru și Ioan le-au răspuns: “Judecați voi înșivă dacă este drept înaintea lui Dumnezeu să vă ascultăm pe voi și nu pe Dumnezeu,

\VS{20}căci noi nu putem să nu spunem ce am văzut și ce am auzit.”


\par }{\PP \VS{21}După ce i-au amenințat și mai mult, i-au lăsat să plece, fără să găsească o cale de a-i pedepsi, din cauza poporului, căci toți slăveau pe Dumnezeu pentru ceea ce se făcuse.

\VS{22}Căci omul asupra căruia s-a făcut această minune de vindecare avea mai mult de patruzeci de ani.


\par }{\PP \VS{23}După ce li s-a dat drumul, au venit în ceata lor și au povestit tot ce le spuseseră preoții cei mai de seamă și bătrânii.

\VS{24}După ce au auzit, și-au înălțat glasul către Dumnezeu într-un glas și au zis: “Doamne, Tu ești Dumnezeu, care ai făcut cerul, pământul, marea și tot ce este în ele;

\VS{25}care, prin gura robului Tău David, ai zis: “Doamne, Tu ești Dumnezeu, care ai făcut cerul, pământul, marea și tot ce este în ele; care, prin gura robului Tău David, ai spus

\par }{\Q 'De ce se înfurie națiunile,

\par }{\QB Și popoarele uneltesc un lucru deșert?


\par }{\Q \VS{26}Împărații pământului iau atitudine,

\par }{\QB și conducătorii complotează împreună,

\par }{\QB împotriva Domnului și împotriva Hristosului Său”.
\FTNT{†}{{\FR 4:26 }{\FT{Hristos (greacă) și Mesia (ebraică) înseamnă amândouă “Unsul”.}}}\FTNT{‡}{{\FR 4:26 }{\FT{{\CHAR{Psalmul 2:1-2}}}}}


\par }{\PP \VS{27}Căci, în adevăr,
\FTNT{§}{{\FR 4:27 }{\FT{TR și NU omit “Isus” și inversează ordinea versetelor 28 și 29.}}}atât Irod, cât și Ponțiu Pilat, cu neamurile și cu poporul lui Israel, s-au adunat împotriva robului Tău cel sfânt, Isus, pe care L-ai uns,

\VS{28}ca să facă tot ce a hotărât mâna Ta și sfatul Tău să se întâmple.

\VS{29}Acum, Doamne, privește la amenințările lor și dă-le slujitorilor tăi să rostească cuvântul tău cu toată îndrăzneala,

\VS{30}în timp ce tu îți întinzi mâna pentru a vindeca și pentru ca semnele și minunile să se facă prin numele sfântului tău slujitor Isus.”


\par }{\PP \VS{31}După ce s-au rugat, s-a cutremurat locul unde erau adunați. Toți au fost umpluți de Duhul Sfânt și au rostit cu îndrăzneală cuvântul lui Dumnezeu.


\par }{\PP \VS{32}Mulțimea celor ce au crezut era de o inimă și de un suflet. Nici unul dintre ei nu pretindea că ceva din lucrurile pe care le poseda îi aparținea, ci aveau toate lucrurile în comun.

\VS{33}Cu mare putere, apostolii dădeau mărturie despre învierea Domnului Isus. Mare har era peste ei toți.

\VS{34}Nu era printre ei nici unul care să ducă lipsă, căci toți cei care erau proprietari de terenuri sau de case le vindeau și aduceau veniturile din lucrurile vândute

\VS{35}și le puneau la picioarele apostolilor; și se făcea împărțirea la fiecare, după cum avea nevoie.


\par }{\PP \VS{36}Iosif, care era numit de apostoli și Barnaba (care înseamnă, în traducere: Fiul încurajării), levit, om din Cipru,

\VS{37}care aveaun câmp, l-a vândut, a adus banii și i-a pus la picioarele apostolilor.



\par }\Chap{5}{\PP \VerseOne{1}Dar un om, numit Anania, împreună cu soția sa Safira, au vândut un bun

\VS{2}și au păstrat o parte din preț, iar soția sa, care știa de aceasta, a adus o parte și a pus-o la picioarele apostolilor.

\VS{3}Dar Petru a zis: “Anania, de ce ți-a umplut Satana inima ca să minți pe Duhul Sfânt și să reții o parte din prețul pământului?

\VS{4}Cât timp l-ai păstrat, nu a rămas oare al tău? După ce a fost vândut, nu a fost în puterea ta? Cum se face că ai conceput acest lucru în inima ta? Nu i-ai mințit pe oameni, ci pe Dumnezeu.”


\par }{\PP \VS{5}Anania, auzind aceste cuvinte, a căzut jos și a murit. O mare frică a cuprins pe toți cei care au auzit aceste lucruri.

\VS{6}Tinerii s-au sculat, l-au înfășurat, l-au dus afară și l-au îngropat.

\VS{7}După vreo trei ore, soția lui, care nu știa ce se întâmplase, a intrat.

\VS{8}Petru i-a răspuns: “Spune-mi dacă ai vândut pământul cu atât de mult.”

\par }{\PP Ea a spus: “Da, pentru atât de mult”.


\par }{\PP \VS{9}Dar Petru a întrebat-o: “Cum de v-ați învoit împreună să ispitiți Duhul Domnului? Iată, picioarele celor care l-au îngropat pe soțul tău sunt la ușă și te vor scoate afară”.


\par }{\PP \VS{10}Ea a căzut îndată la picioarele lui și a murit. Tinerii au intrat și au găsit-o moartă; au dus-o afară și au îngropat-o lângă soțul ei.

\VS{11}O mare frică a cuprins toată adunarea și pe toți cei care au auzit aceste lucruri.


\par }{\PP \VS{12}Prin mâinile apostolilor se făceau multe semne și minuni în popor. Toți erau de acord în pridvorul lui Solomon.

\VS{13}Nici unul dintre ceilalți nu îndrăznea să li se alăture; totuși, poporul îi cinstea.

\VS{14}Tot mai mulți credincioși s-au adăugat la Domnul, mulțimi de bărbați și de femei.

\VS{15}Chiar și pe bolnavi îi scoteau în stradă și îi așezau pe paturi și pe saltele, pentru ca, atunci când Petru trecea, măcar umbra lui să-i umbrească pe unii dintre ei.

\VS{16}De asemenea, s-a adunat o mulțime din cetățile din jurul Ierusalimului, aducând bolnavi și pe cei chinuiți de duhuri necurate; și toți au fost vindecați.


\par }{\PP \VS{17}Dar marele preot și toți cei ce erau cu el, adică secta saducheilor, s-au sculat în picioare și s-au umplut de gelozie

\VS{18}și au pus mâna pe apostoli, apoi i-au pus în temniță publică.

\VS{19}Dar un înger al Domnului a deschis noaptea ușile închisorii, i-a scos afară și le-a spus:

\VS{20}“Mergeți și stați în picioare și vorbiți poporului în templu toate cuvintele acestei vieți”.


\par }{\PP \VS{21}După ce au auzit acestea, au intrat în Templu în zori de ziuă și învățau. Dar marele preot și cei care erau cu el au venit și au convocat consiliul, împreună cu tot senatul copiilor lui Israel, și au trimis la închisoare ca să fie aduși.

\VS{22}Dar ofițerii care au venit nu i-au găsit în închisoare. S-au întors și au raportat:

\VS{23}“Am găsit închisoarea închisă și încuiată, iar gardienii stăteau în fața ușilor, dar când le-am deschis, nu am găsit pe nimeni înăuntru!”


\par }{\PP \VS{24}Când au auzit cuvintele acestea, marele preot, căpetenia templului și preoții cei mai de seamă au fost foarte nedumeriți de ele și de ceea ce putea să se întâmple.

\VS{25}Cineva a venit și le-a spus: “Iată că oamenii pe care i-ați pus în închisoare sunt în templu, stând în picioare și învățând poporul.”

\VS{26}Atunci căpitanul s-a dus cu ofițerii și i-a adus fără violență, căci se temeau că poporul îi va ucide cu pietre.


\par }{\PP \VS{27}După ce i-au adus, i-au pus în fața consiliului. Marele preot i-a interpelat și le-a zis:

\VS{28}“Nu v-am poruncit noi cu strictețe să nu învățați în acest nume? Iată că ați umplut Ierusalimul cu învățătura voastră și intenționați să aduceți asupra noastră sângele acestui om.”


\par }{\PP \VS{29}Dar Petru și apostolii au răspuns: “Trebuie să ascultăm de Dumnezeu mai mult decât de oameni.

\VS{30}Dumnezeul părinților noștri a înviat pe Isus, pe care voi L-ați omorât, spânzurându-L pe un lemn.

\VS{31}Dumnezeu L-a înălțat cu dreapta Sa ca să fie Prinț și Mântuitor, ca să dea pocăință lui Israel și iertare de păcate.

\VS{32}Noi suntem martorii Lui despre aceste lucruri; la fel și Duhul Sfânt, pe care Dumnezeu l-a dat celor care Îl ascultă.”


\par }{\PP \VS{33}Dar ei, auzind aceasta, au fost înduioșați și au hotărât să-i omoare.

\VS{34}Dar unul dintre ei s-a ridicat în consiliu, un fariseu numit Gamaliel, învățător al legii, cinstit de tot poporul, și a poruncit să-i scoată pe apostoli pentru puțină vreme.

\VS{35}El le-a zis: “Bărbați israeliți, fiți atenți la acești oameni, la ceea ce aveți de gând să faceți.

\VS{36}Pentru că înainte de aceste zile s-a ridicat Theudas, dându-se drept cineva, căruia i s-au alăturat un număr de oameni, cam patru sute. El a fost ucis; și toți, câți l-au ascultat, s-au risipit și au rămas fără nimic.

\VS{37}După acesta, Iuda din Galileea s-a ridicat în zilele înrobirii și a atras după el niște oameni. Și el a pierit și el, și toți, câți îl ascultau, au fost împrăștiați.

\VS{38}Acum vă spun să vă îndepărtați de acești oameni și să-i lăsați în pace. Căci, dacă acest sfat sau această lucrare este de la oameni, va fi răsturnată.

\VS{39}Dar dacă este de la Dumnezeu, nu veți putea să o răsturnați și veți fi găsiți chiar că luptați împotriva lui Dumnezeu!”


\par }{\PP \VS{40}Ei au fost de acord cu el. Chemându-i pe apostoli, i-au bătut, le-au poruncit să nu mai vorbească în numele lui Isus și i-au lăsat să plece.

\VS{41}Au plecat deci din fața consiliului, bucurându-se că au fost socotiți vrednici să sufere dezonoare pentru numele lui Isus.


\par }{\PP \VS{42}În fiecare zi, în templu și în casele lor, nu încetau să învețe și să propovăduiască pe Iisus Hristosul.



\par }\Chap{6}{\PP \VerseOne{1}În zilele acelea, când numărul ucenicilor se înmulțea, s-a ridicat o plângere din partea elinilor împotriva evreilor, pentru că văduvele lor erau neglijate în slujba zilnică.

\VS{2}Cei doisprezece au convocat mulțimea discipolilor și au spus: “Nu se cuvine ca noi să părăsim cuvântul lui Dumnezeu și să servim la mese.

\VS{3}De aceea, fraților, alegeți dintre voi șapte bărbați de bună reputație, plini de Duhul Sfânt și de înțelepciune, pe care să îi numim în fruntea acestei afaceri.

\VS{4}Dar noi vom continua cu stăruință în rugăciune și în slujirea cuvântului.”


\par }{\PP \VS{5}Aceste cuvinte au plăcut întregii mulțimi. Ei au ales pe Ștefan, un om plin de credință și de Duhul Sfânt, pe Filip, Prohor, Nicanor, Timon, Parmenas și Nicolae, un prozelit din Antiohia,

\VS{6}pe care i-au pus înaintea apostolilor. După ce s-au rugat, aceștia și-au pus mâinile peste ei.


\par }{\PP \VS{7}Cuvântul lui Dumnezeu creștea și numărul ucenicilor se înmulțea foarte mult în Ierusalim. O mare mulțime de preoți se supuneau credinței.


\par }{\PP \VS{8}Ștefan, plin de credință și de putere, făcea minuni și semne mari în popor.

\VS{9}Dar unii dintre cei din sinagoga numită “Libertinii”, dintre cireni, dintre alexandrini și dintre cei din Cilicia și din Asia s-au sculat, certându-se cu Ștefan.

\VS{10}Ei nu au putut să reziste înțelepciunii și Duhului prin care vorbea el.

\VS{11}Atunci au determinat pe ascuns pe oameni să spună: “L-am auzit spunând cuvinte blasfemiatoare împotriva lui Moise și a lui Dumnezeu”.

\VS{12}Au stârnit poporul, bătrânii și cărturarii, au venit împotriva lui și l-au prins, apoi l-au adus în consiliu,

\VS{13}și au pus martori falși care spuneau: “Omul acesta nu încetează să rostească cuvinte blasfemiatoare împotriva acestui loc sfânt și a Legii.

\VS{14}Căci l-am auzit spunând că acest Isus din Nazaret va distruge acest loc și va schimba obiceiurile pe care ni le-a dat Moise.”

\VS{15}Toți cei care stăteau în consiliu, fixându-și ochii asupra lui, i-au văzut fața ca pe cea a unui înger.



\par }\Chap{7}{\PP \VerseOne{1}Marele preot a zis: “Așa stau lucrurile?”


\par }{\PP \VS{2}El a zis: “Frați și părinți, ascultați. Dumnezeul slavei s-a arătat tatălui nostru Avraam, pe când era în Mesopotamia, înainte de a locui în Haran,

\VS{3}și i-a spus: “Ieși din țara ta și de lângă rudele tale și vino într-o țară pe care ți-o voi arăta.

\VS{4}Atunci a ieșit din țara caldeenilor și a locuit în Haran. De acolo, după ce tatăl său a murit, Dumnezeu l-a mutat în acest pământ în care locuiți acum.

\VS{5}Nu i-a dat nicio moștenire în ea, nu, nici măcar ca să pună piciorul pe ea. I-a promis că i-o va da ca moștenire și urmașilor lui după el, când el nu avea încă niciun copil.

\VS{6}Dumnezeu a vorbit astfel: că urmașii lui vor trăi ca niște străini într-o țară străină și că vor fi înrobiți și maltratați timp de patru sute de ani.

\VS{7}“Voi judeca națiunea la care vor fi în robie”, a spus Dumnezeu, “și după aceea vor ieși și îmi vor sluji în acest loc.

\VS{8}Și i-a dat legământul circumciziei. Astfel, Avraam a devenit tatăl lui Isaac și l-a circumcis în ziua a opta. Isaac a devenit tatăl lui Iacob, iar Iacob a devenit tatăl celor doisprezece patriarhi.


\par }{\PP \VS{9}“Patriarhii au fost geloși pe Iosif și l-au vândut în Egipt. Dumnezeu a fost cu el

\VS{10}și l-a izbăvit din toate necazurile sale și i-a dat favoare și înțelepciune înaintea lui Faraon, regele Egiptului. L-a pus guvernator peste Egipt și peste toată casa lui.

\VS{11}Și a venit o foamete peste toată țara Egiptului și Canaanului și o mare nenorocire. Părinții noștri nu au găsit hrană.

\VS{12}Dar când a auzit Iacov că în Egipt era grâu, a trimis prima dată pe părinții noștri.

\VS{13}A doua oară, Iosif s-a făcut cunoscut fraților săi, iar familia lui Iosif a fost descoperită lui Faraon.

\VS{14}Iosif a trimis și a chemat pe tatăl său Iacov și pe toate rudele lui, șaptezeci și cinci de suflete.

\VS{15}Iacov a coborât în Egipt și a murit, el însuși și părinții noștri;

\VS{16}au fost aduși înapoi la Sihem și au fost puși în mormântul pe care Avraam îl cumpărase pe un preț în argint de la fiii lui Hamor din Sihem.


\par }{\PP \VS{17}Dar, când s-a apropiat vremea făgăduinței pe care Dumnezeu o făgăduise lui Avraam, poporul a crescut și s-a înmulțit în Egipt,

\VS{18}până când s-a ridicat un alt rege, care nu-l cunoștea pe Iosif.

\VS{19}Acesta a profitat de neamul nostru și i-a maltratat pe părinții noștri, obligându-i să își abandoneze copiii, ca să nu rămână în viață.

\VS{20}În acel timp s-a născut Moise, care era extrem de frumos pentru Dumnezeu. El a fost hrănit trei luni în casa tatălui său.

\VS{21}După ce a fost abandonat, fiica lui Faraon l-a luat în brațe și l-a crescut ca pe propriul ei fiu.

\VS{22}Moise a fost instruit în toată înțelepciunea egiptenilor. A fost puternic în cuvintele și în faptele sale.

\VS{23}Dar când a împlinit patruzeci de ani, i-a venit în inimă să viziteze pe frații săi, copiii lui Israel.

\VS{24}Văzând că unul dintre ei suferă nedreptate, l-a apărat și l-a răzbunat pe cel asuprit, lovindu-l pe egiptean.

\VS{25}El presupunea că frații săi au înțeles că Dumnezeu, prin mâna lui, le dădea izbăvirea; dar ei nu au înțeles.


\par }{\PP \VS{26}A doua zi, când se luptau, li s-a arătat și i-a îndemnat să se împace din nou, zicând: “Domnilor, voi sunteți frați. De ce vă faceți rău unul altuia?”

\VS{27}Dar cel care îi făcea rău aproapelui său îl respingea, spunând: “Cine te-a pus conducător și judecător peste noi?

\VS{28}Vrei să mă omori pe mine, cum ai omorât ieri pe egiptean?”

\VS{29}La acest cuvânt, Moise a fugit și a devenit străin în țara lui Madian, unde a devenit tatăl a doi fii.


\par }{\PP \VS{30}“După ce s-au împlinit patruzeci de ani, un înger al Domnului i s-a arătat în pustiul muntelui Sinai, într-o flacără de foc, într-un tufiș.

\VS{31}Când a văzut-o, Moise s-a mirat de această priveliște. Pe când se apropia să vadă, vocea Domnului i s-a adresat:

\VS{32}“Eu sunt Dumnezeul părinților tăi: Dumnezeul lui Avraam, Dumnezeul lui Isaac și Dumnezeul lui Iacov. Moise a tremurat și nu a îndrăznit să privească.

\VS{33}Domnul i-a zis: “Scoate-ți sandalele, căci locul unde stai este pământ sfânt.

\VS{34}Am văzut cu adevărat necazul poporului Meu care este în Egipt și i-am auzit gemetele. M-am coborât să-i eliberez. Acum vino, te voi trimite în Egipt”.


\par }{\PP \VS{35}Acest Moise, pe care ei l-au refuzat, zicând: “Cine te-a făcut pe tine conducător și judecător?” Dumnezeu l-a trimis ca conducător și eliberator prin mâna îngerului care i s-a arătat în tufiș.

\VS{36}Acest om i-a condus, după ce a făcut minuni și semne în Egipt, în Marea Roșie și în pustiu timp de patruzeci de ani.

\VS{37}Acesta este acel Moise care le-a spus copiilor lui Israel: “Domnul Dumnezeul nostru vă va ridica un profet dintre frații voștri, ca mine”.

\VS{38}Acesta este cel care a fost în adunare în pustiu cu îngerul care i-a vorbit pe muntele Sinai și cu părinții noștri, care a primit revelații vii pe care să ni le dea nouă,

\VS{39}căruia părinții noștri nu au vrut să fie ascultători, ci l-au respins și s-au întors cu inima în Egipt,

\VS{40}spunând lui Aaron: 'Fă-ne dumnezei care să meargă înaintea noastră, căci în ceea ce-l privește pe acest Moise care ne-a scos din țara Egiptului, nu știm ce s-a întâmplat cu el'.

\VS{41}În zilele acelea au făcut un vițel, au adus jertfe idolului și s-au bucurat de lucrările mâinilor lor.

\VS{42}Dar Dumnezeu s-a întors de la ei și i-a predat să slujească armatei cerului, după cum este scris în cartea profeților,

\par }{\Q “Mi-ai oferit animale ucise și sacrificii?

\par }{\QB Patruzeci de ani în pustie, casă a lui Israel?


\par }{\Q \VS{43}Tu ai ridicat cortul lui Moloh,

\par }{\QB steaua zeului tău, Rephan,

\par }{\Q figurile pe care le-ai făcut să se închine,

\par }{\QB așa că te voi duce dincolo de Babilon.


\par }{\PP \VS{44}“Părinții noștri au avut cortul mărturiei în pustie, așa cum îi poruncise Cel ce a vorbit lui Moise să îl facă după modelul pe care îl văzuse;

\VS{45}și pe care, la rândul lor, părinții noștri l-au adus cu Iosua, când au intrat în stăpânirea neamurilor pe care Dumnezeu le-a alungat dinaintea părinților noștri până în zilele lui David,

\VS{46}care au găsit favoare înaintea lui Dumnezeu și au cerut să găsească o locuință pentru Dumnezeul lui Iacov.

\VS{47}Dar Solomon i-a zidit o casă.

\VS{48}Totuși, Cel Preaînalt nu locuiește în temple făcute de mâini, cum spune profetul,


\par }{\Q \VS{49}“Cerul este tronul meu,

\par }{\QB și pământul un scăunel pentru picioarele mele.

\par }{\Q Ce fel de casă îmi vei construi?”, zice Domnul.

\par }{\QB “Sau care este locul meu de odihnă?


\par }{\Q \VS{50}Nu cumva mâna mea a făcut toate aceste lucruri?


\par }{\PP \VS{51}“Voi, cei cu gâtul împietrit și netăiați împrejur la inimă și la urechi, vă împotriviți mereu Duhului Sfânt! Cum au făcut părinții voștri, așa faceți și voi.

\VS{52}Pe care dintre profeți nu i-au prigonit părinții voștri? Ei i-au ucis pe cei care au prezis venirea Celui Drept, dintre care voi ați devenit acum trădători și ucigași.

\VS{53}Ați primit Legea așa cum a fost rânduită de îngeri și nu ați respectat-o!”


\par }{\PP \VS{54}Când au auzit acestea, li s-a tăiat inima și au scrâșnit din dinți împotriva Lui.

\VS{55}Dar el, fiind plin de Duhul Sfânt, a privit cu stăruință spre cer și a văzut gloria lui Dumnezeu și pe Isus stând la dreapta lui Dumnezeu,

\VS{56}și a zis: “Iată, văd cerurile deschise și pe Fiul Omului stând la dreapta lui Dumnezeu!”


\par }{\PP \VS{57}Dar ei au strigat cu glas tare și și-au astupat urechile, apoi s-au năpustit la unison asupra lui.

\VS{58}L-au aruncat afară din cetate și l-au ucis cu pietre. Martorii și-au așezat hainele la picioarele unui tânăr pe nume Saul.

\VS{59}L-au ucis cu pietre pe Ștefan, pe când striga: “Doamne Iisuse, primește duhul meu!”

\VS{60}El a îngenuncheat și a strigat cu glas tare: “Doamne, nu le reține acest păcat!” După ce a spus aceasta, a adormit.



\par }\Chap{8}{\PP \VerseOne{1}Saul consimțea la moartea sa. O mare persecuție s-a ridicat împotriva adunării care era în Ierusalim în acea zi. Toți erau împrăștiați prin regiunile Iudeii și Samariei, cu excepția apostolilor.

\VS{2}Oameni evlavioși l-au îngropat pe Ștefan și l-au plâns mult.

\VS{3}Dar Saul a devastat Adunarea, a intrat în fiecare casă și a târât în închisoare atât bărbați cât și femei.

\VS{4}De aceea, cei care erau împrăștiați prin împrejurimi mergeau de colo-colo, propovăduind cuvântul.

\VS{5}Filip s-a coborât în cetatea Samaria și le-a vestit pe Hristos.

\VS{6}Mulțimile ascultau cu un singur glas cele spuse de Filip, când auzeau și vedeau semnele pe care le făcea.

\VS{7}Căci duhuri necurate ieșeau din mulți dintre cei care le aveau. Ieșeau, strigând cu glas tare. Mulți dintre cei care fuseseră paralizați și șchiopi au fost vindecați.

\VS{8}A fost o mare bucurie în cetatea aceea.


\par }{\PP \VS{9}Dar era un om, pe nume Simon, care făcea vrăjitorii în cetate și care uimea pe locuitorii Samariei, dându-se drept un mare om.

\VS{10}Toți ascultau, de la cel mai mic până la cel mai mare, și ziceau: “Acesta este o mare putere a lui Dumnezeu.”

\VS{11}Îl ascultau pentru că de mult timp îi uimise cu vrăjitoriile sale.

\VS{12}Dar, când au crezut că Filip propovăduia vestea bună despre Împărăția lui Dumnezeu și despre numele lui Isus Hristos, s-au botezat, atât bărbații, cât și femeile.

\VS{13}Și Simon însuși a crezut. Fiind botezat, a continuat să meargă cu Filip. Văzând că se întâmplă semne și minuni mari, a rămas uimit.


\par }{\PP \VS{14}Apostolii care erau la Ierusalim, auzind că Samaria a primit cuvântul lui Dumnezeu, au trimis la ei pe Petru și pe Ioan,

\VS{15}care, după ce s-au coborât, s-au rugat pentru ei, ca să primească Duhul Sfânt,

\VS{16}pentru că încă nu se pogorâse peste niciunul dintre ei. Ei fuseseră doar botezați în numele lui Isus Cristos.

\VS{17}Atunci și-au pus mâinile peste ei și au primit Duhul Sfânt.

\VS{18}Simon, văzând că Duhul Sfânt se dădea prin punerea mâinilor apostolilor, le-a oferit bani,

\VS{19}zicând: “Dați-mi și mie această putere, ca oricine îmi voi pune mâinile să primească Duhul Sfânt”.

\VS{20}Dar Petru i-a spus: “Să piară argintul tău împreună cu tine, pentru că ai crezut că poți obține darul lui Dumnezeu cu bani!

\VS{21}Tu nu ai nici parte, nici sorți de izbândă în această chestiune, pentru că inima ta nu este dreaptă înaintea lui Dumnezeu.

\VS{22}Pocăiește-te, așadar, de aceasta, de răutatea ta, și întreabă-l pe Dumnezeu dacă nu cumva gândul inimii tale poate să-ți fie iertat.

\VS{23}Căci văd că ești în otrava amărăciunii și în robia nelegiuirii.”


\par }{\PP \VS{24}Simon a răspuns: “Roagă-te pentru mine Domnului, ca să nu mi se întâmple nimic din cele ce ai spus.”


\par }{\PP \VS{25}După ce au mărturisit și au rostit cuvântul Domnului, s-au întors la Ierusalim și au propovăduit vestea cea bună în multe sate ale samaritenilor.


\par }{\PP \VS{26}Atunci un înger al Domnului a vorbit lui Filip și i-a zis: “Scoală-te și mergi spre sud, pe drumul care coboară de la Ierusalim la Gaza. Acesta este un deșert”.


\par }{\PP \VS{27}S-a sculat și s-a dus; și iată că era un bărbat din Etiopia, un eunuc cu mare putere sub Candace, regina etiopienilor, care era stăpână pe toată comoara ei, și care venise la Ierusalim să se închine.

\VS{28}El se întorcea, ședea în carul său și citea proorocul Isaia.


\par }{\PP \VS{29}Duhul Sfânt a zis lui Filip: “Apropie-te și ia-te de carul acesta.”


\par }{\PP \VS{30}Filip a alergat la el, l-a auzit citind pe Isaia proorocul și i-a zis: “Înțelegi tu ce citești?”


\par }{\PP \VS{31}El a zis: “Cum aș putea, dacă nu-mi explică cineva?” L-a rugat pe Filip să urce și să stea cu el.

\VS{32}Și pasajul din Scriptură pe care îl citea el era acesta,

\par }{\Q “A fost dus ca o oaie la abator.

\par }{\QB Ca un miel în fața celui ce-l tunde tace,

\par }{\QB ca să nu deschidă gura.


\par }{\Q \VS{33}În umilirea Lui, a fost luată judecata Lui.

\par }{\QB Cine va declara generația Lui?

\par }{\QB Căci viața lui este luată de pe pământ.”


\par }{\PP \VS{34}Eunucul a răspuns lui Filip: “Despre cine vorbește proorocul acesta? Despre el însuși sau despre altcineva?”


\par }{\PP \VS{35}Filip a deschis gura și, pornind de la Scriptura aceasta, i-a propovăduit despre Isus.

\VS{36}Pe când mergeau pe drum, au ajuns la o apă; și eunucul a zis: “Iată, aici este apă. Ce mă împiedică să fiu botezat?”


\par }{\PP \VS{37}\FTNT{*}{{\FR 8:37 }{\FT{TR adaugă
}}{\FQA{Filip a spus: „Dacă crezi din toată inima, poți.” El a răspuns: „Cred că Isus Hristos este Fiul lui Dumnezeu”.}}}

\VS{38}Și a poruncit să se oprească carul și s-au coborât amândoi în apă, atât Filip cât și eunucul, și l-a botezat.


\par }{\PP \VS{39}După ce au ieșit din apă, Duhul Domnului a răpit pe Filip, și eunucul nu l-a mai văzut, căci a plecat bucuros.

\VS{40}Dar Filip a fost găsit la Azotus. Trecând pe acolo, a propovăduit Vestea cea Bună în toate orașele, până când a ajuns la Cezareea.



\par }\Chap{9}{\PP \VerseOne{1}Dar Saul, care continua să arunce amenințări și să ucidă pe ucenicii Domnului, s-a dus la marele preot

\VS{2}și a cerut de la el scrisori către sinagogile din Damasc, ca, dacă va găsi pe vreunul din cei de pe Cale, fie bărbați, fie femei, să îl aducă legat la Ierusalim.

\VS{3}În timp ce călătorea, s-a apropiat de Damasc și, deodată, o lumină din cer a strălucit în jurul lui.

\VS{4}A căzut la pământ și a auzit un glas care îi spunea:
{\WJ{“Saul, Saul, de ce mă prigonești?”.
}}


\par }{\PP \VS{5}El a zis: “Cine ești Tu, Doamne?”

\par }{\PP Domnul a spus: “{\WJ{Eu sunt Isus, pe care voi Îl prigoniți.
}}

\VS{6}{\WJ{Dar, sculându-te, intră în cetate și atunci ți se va spune ce trebuie să faci.”
}}


\par }{\PP \VS{7}Oamenii care călătoreau cu el au rămas muți, auzind zgomotul, dar nu vedeau pe nimeni.

\VS{8}Saul s-a ridicat de la pământ și, când și-a deschis ochii, n-a văzut pe nimeni. L-au luat de mână și l-au dus în Damasc.

\VS{9}A stat fără vedere timp de trei zile și nu a mâncat și nu a băut.


\par }{\PP \VS{10}Și era în Damasc un ucenic, numit Anania. Domnul i-a spus într-o viziune:
{\WJ{“Anania!”
}}

\par }{\PP El a spus: “Iată, sunt eu, Doamne”.


\par }{\PP \VS{11}Domnul i-a zis:
{\WJ{“Scoală-te și du-te pe strada care se cheamă Drept și caută în casa lui Iuda pe unul numit Saul, un om din Tars. Căci iată că el se roagă
}}

\VS{12}și a
{\WJ{văzut în viziune pe un om numit Anania intrând și punându-și mâinile peste el, ca să-și recapete vederea.”
}}


\par }{\PP \VS{13}Și Anania a răspuns: “Doamne, am auzit de la mulți despre omul acesta, cât rău a făcut sfinților Tăi din Ierusalim.

\VS{14}Aici are putere de la preoții cei mai de seamă să lege pe toți cei ce cheamă Numele Tău.”


\par }{\PP \VS{15}Dar Domnul i-a zis: “{\WJ{Du-te, căci el este vasul Meu ales ca să poarte numele Meu înaintea neamurilor, a împăraților și a copiilor lui Israel.
}}

\VS{16}{\WJ{Căci îi voi arăta câte va trebui să sufere pentru Numele Meu.”
}}


\par }{\PP \VS{17}Anania a plecat și a intrat în casă. Și, punându-și mâinile peste el, i-a zis: “Frate Saul, Domnul, care ți s-a arătat pe drumul pe care ai venit, m-a trimis ca să-ți recapeți vederea și să fii umplut de Duhul Sfânt.”

\VS{18}Imediat, ceva ca niște solzi au căzut de pe ochii lui și și-a recăpătat vederea. S-a sculat și a fost botezat.

\VS{19}A luat mâncare și s-a întărit.

\par }{\PP Saul a stat câteva zile cu ucenicii care se aflau la Damasc.

\VS{20}Îndată, în sinagogi, a propovăduit pe Hristos, că El este Fiul lui Dumnezeu.

\VS{21}Toți cei care îl auzeau erau uimiți și ziceau: “Nu cumva acesta este cel care, la Ierusalim, a făcut ravagii printre cei care invocau acest nume? Și venise aici cu intenția de a-i aduce legați în fața preoților de seamă!”.


\par }{\PP \VS{22}Dar Saul, care se întărea și mai mult, a zăpăcit pe iudeii care locuiau în Damasc și a dovedit că acesta este Hristosul.

\VS{23}După ce s-au împlinit multe zile, iudeii au uneltit împreună ca să-l ucidă,

\VS{24}dar complotul lor a fost cunoscut de Saul. Ei pândeau porțile zi și noapte ca să-l ucidă,

\VS{25}dar discipolii lui l-au luat noaptea și l-au coborât prin zid, coborându-l într-un coș.


\par }{\PP \VS{26}După ce a ajuns la Ierusalim, Saul a încercat să se alăture ucenicilor, dar toți se temeau de el, nevenindu-le să creadă că este ucenic.

\VS{27}Dar Barnaba l-a luat și l-a dus la apostoli și le-a povestit cum Îl văzuse pe Domnul pe drum, cum îi vorbise și cum la Damasc predicase cu îndrăzneală în numele lui Isus.

\VS{28}A intrat cu ei în Ierusalim,

\VS{29}și propovăduia cu îndrăzneală în Numele Domnului Isus. El vorbea și se certa cu elenii,
\FTNT{*}{{\FR 9:29 }{\FT{Eleniștii erau evrei care foloseau limba și cultura greacă.}}}dar aceștia căutau să-l ucidă.

\VS{30}Când au aflat frații, l-au coborât în Cezareea și l-au trimis la Tars.


\par }{\PP \VS{31}Și adunările din toată Iudeea, din Galileea și din Samaria au avut pace și s-au întărit. Ele s-au înmulțit, umblând în frica Domnului și în mângâierea Duhului Sfânt.


\par }{\PP \VS{32}Petru a străbătut toate părțile acelea, și s-a pogorât și la sfinții care locuiau în Lida.

\VS{33}Acolo a găsit un om numit Enea, care era țintuit la pat de opt ani, pentru că era paralizat.

\VS{34}Petru i-a zis: “Enea, Isus Hristos te vindecă. Ridică-te și fă-ți patul!” Imediat, el s-a sculat.

\VS{35}Toți cei care locuiau la Lida și în Saron l-au văzut și s-au întors la Domnul.


\par }{\PP \VS{36}În Iope era la Iope o ucenică cu numele Tabita, care în traducere înseamnă Dorcas.\FTNT{†}{{\FR 9:36 }{\FT{“Dorcas” în greacă înseamnă “gazelă”.}}} Această femeie era plină de fapte bune și de acte de milostenie pe care le făcea.

\VS{37}În zilele acelea, ea s-a îmbolnăvit și a murit. După ce au spălat-o, au așezat-o într-o cameră de sus.

\VS{38}Cum Lida era aproape de Iope, ucenicii, auzind că Petru era acolo, au trimis doi bărbați\FTNT{‡}{{\FR 9:38 }{\FT{Citire din NU, TR; MT omite “doi bărbați”}}} la el, rugându-l să nu întârzie să vină la ei.

\VS{39}Petru s-a ridicat și a plecat cu ei. După ce a venit, l-au dus în camera de sus. Toate văduvele stăteau lângă el, plângând și arătând tunicile și celelalte haine pe care Dorcas le făcuse cât timp fusese cu ele.

\VS{40}Petru le-a trimis pe toate afară, a îngenuncheat și s-a rugat. Întorcându-se spre trup, a zis: “Tabita, ridică-te!”. Ea a deschis ochii și, când l-a văzut pe Petru, s-a așezat în picioare.

\VS{41}El i-a dat mâna și a ridicat-o în picioare. Chemând sfinții și văduvele, a prezentat-o în viață.

\VS{42}Acest lucru s-a aflat în toată Ioppa și mulți au crezut în Domnul.

\VS{43}A stat multe zile în Iope la un tăbăcar cu numele Simon.



\par }\Chap{10}{\PP \VerseOne{1}În Cezareea era un om, Corneliu, centurion în ceea ce se numea regimentul italian,

\VS{2}un om evlavios și temător de Dumnezeu în toată casa lui, care dădea cu generozitate poporului daruri pentru cei nevoiași și se ruga mereu lui Dumnezeu.

\VS{3}Pe la ora nouă a zilei,
\FTNT{*}{{\FR 10:3 }{\FT{3:00 p.m.}}}a văzut clar într-o viziune un înger al lui Dumnezeu care a venit la el și i-a spus: “Corneliu!”.


\par }{\PP \VS{4}El, care, fixându-și ochii asupra Lui și înfricoșându-se, a zis: “Ce este, Doamne?”

\par }{\PP El i-a spus: “Rugăciunile tale și darurile tale pentru cei nevoiași s-au înălțat ca o amintire înaintea lui Dumnezeu.

\VS{5}Trimite acum oameni la Iope și adu-l pe Simon, care se mai numește și Petru.

\VS{6}El locuiește la un tăbăcar pe nume Simon, a cărui casă se află pe malul mării.
\FTNT{†}{{\FR 10:6 }{\FT{TR adaugă: “Acesta vă va spune ce trebuie să faceți”.}}}


\par }{\PP \VS{7}După ce a plecat îngerul care-i vorbise, Corneliu a chemat doi dintre slujitorii săi de casă și un ostaș credincios dintre cei care-l îngrijeau în permanență.

\VS{8}După ce le-a explicat totul, i-a trimis la Iope.


\par }{\PP \VS{9}A doua zi, când au pornit în călătorie și s-au apropiat de cetate, Petru s-a suit pe acoperișul casei să se roage, pe la amiază.

\VS{10}I s-a făcut foame și a vrut să mănânce, dar, în timp ce ei se pregăteau, a căzut în transă.

\VS{11}A văzut cerul deschis și un anumit recipient care se cobora la el, ca un cearșaf mare lăsat de patru colțuri pe pământ,

\VS{12}în care se aflau tot felul de animale cu patru picioare de pe pământ, animale sălbatice, reptile și păsări ale cerului.

\VS{13}Un glas a venit la el:
{\WJ{“Scoală-te, Petru, ucide și mănâncă!”
}}


\par }{\PP \VS{14}Dar Petru a zis: “Nu, Doamne, căci eu n-am mâncat niciodată nimic din ce este de rând sau necurat.”


\par }{\PP \VS{15}Și a doua oară, un glas a venit din nou la el: “{\WJ{Ceea ce a curățat Dumnezeu, să nu numești necurat.”
}}

\VS{16}A făcut aceasta de trei ori și îndată lucrul acesta a fost primit în cer.


\par }{\PP \VS{17}Și pe când Petru era foarte nedumerit în sinea lui cu privire la semnificația viziunii pe care o avusese, iată că oamenii trimiși de Corneliu, după ce au întrebat de casa lui Simon, au stat în fața porții

\VS{18}și au strigat și au întrebat dacă Simon, care se mai numea și Petru, era găzduit acolo.

\VS{19}În timp ce Petru reflecta asupra viziunii, Duhul i-a spus: “Iată, trei\FTNT{‡}{{\FR 10:19 }{\FT{Citirea din TR și NU. MT omite “trei”}}} bărbați te caută.

\VS{20}Dar ridică-te, coboară-te și mergi cu ei, fără să te îndoiești de nimic, pentru că Eu i-am trimis”.


\par }{\PP \VS{21}Petru s-a coborât la oamenii aceia și le-a zis: “Iată, eu sunt cel pe care-l căutați. De ce ați venit?”


\par }{\PP \VS{22}Ei au zis: “Corneliu, centurionul, un om drept și temător de Dumnezeu, bine pomenit de tot neamul iudeilor, a fost îndemnat de un înger sfânt să vă cheme în casa lui și să asculte ce veți spune.”

\VS{23}Așa că i-a chemat și le-a pus la dispoziție un loc de cazare.

\par }{\PP A doua zi, Petru s-a sculat și a ieșit cu ei; și câțiva dintre frații din Iope l-au însoțit.

\VS{24}A doua zi, au intrat în Cezareea. Corneliu îi aștepta, după ce-și adunase rudele și apropiații.

\VS{25}Când a intrat Petru, Corneliu l-a întâmpinat, a căzut la picioarele lui și i s-a închinat.

\VS{26}Dar Petru l-a ridicat și i-a zis: “Ridică-te! Și eu însumi sunt un om”.

\VS{27}În timp ce vorbea cu el, a intrat și a găsit pe mulți adunați.

\VS{28}El le-a zis: “Voi înșivă știți cum este un lucru nelegiuit ca un iudeu să se unească sau să vină la unul de alt neam, dar Dumnezeu mi-a arătat că nu trebuie să numesc pe nimeni profan sau necurat.

\VS{29}De aceea am venit și eu fără să mă plâng când am fost chemat. De aceea întreb: “De ce ați trimis după mine?”.


\par }{\PP \VS{30}Corneliu a zis: “Acum patru zile am postit până la ceasul acesta; și la ceasul al nouălea,
\FTNT{§}{{\FR 10:30 }{\FT{15:00 p.m.}}}m-am rugat în casa mea și iată că un om îmbrăcat în haine strălucitoare

\VS{31}a stat înaintea mea și a zis: “Corneliu, rugăciunea ta a fost ascultată și darurile tale pentru cei nevoiași au fost luate în seamă înaintea lui Dumnezeu.

\VS{32}Trimite deci la Iope și cheamă-l pe Simon, care se mai numește și Petru. El locuiește în casa unui tăbăcar numit Simon, pe malul mării. Când va veni, va vorbi cu voi”.

\VS{33}De aceea am trimis îndată la tine și a fost bine că ai venit. De aceea, acum suntem cu toții prezenți aici, în fața lui Dumnezeu, ca să auzim tot ce ți-a poruncit Dumnezeu.”


\par }{\PP \VS{34}Petru a deschis gura și a zis: “Cu adevărat, am văzut că Dumnezeu nu face favoruri,

\VS{35}ci, în orice neam, cel ce se teme de El și face dreptate este plăcut Lui.

\VS{36}Cuvântul pe care l-a trimis copiilor lui Israel, propovăduind vestea bună a păcii prin Isus Cristos — el este Domnul tuturor—

\VS{37}voi înșivă știți ce s-a întâmplat, care a fost vestit în toată Iudeea, începând din Galileea, după botezul pe care l-a predicat Ioan;

\VS{38}cum Dumnezeu l-a uns cu Duh Sfânt și cu putere pe Isus din Nazaret, care umbla făcând bine și vindecând pe toți cei asupriți de diavol, pentru că Dumnezeu era cu el.

\VS{39}Noi suntem martori la tot ce a făcut El, atât în țara iudeilor, cât și în Ierusalim, pe care l-au și\FTNT{*}{{\FR 10:39 }{\FT{TR omite “de asemenea”}}} ucis, spânzurându-l pe un lemn.

\VS{40}Dumnezeu l-a înviat a treia zi și l-a dat să se arate,

\VS{41}nu întregului popor, ci martorilor care au fost aleși mai dinainte de Dumnezeu, nouă, care am mâncat și am băut cu el după ce a înviat din morți.

\VS{42}Ne-a poruncit să propovăduim poporului și să mărturisim că Acesta este Cel desemnat de Dumnezeu ca Judecător al celor vii și al celor morți.

\VS{43}Toți profeții mărturisesc despre el că, prin numele lui, oricine crede în el va primi iertarea păcatelor.”


\par }{\PP \VS{44}Pe când Petru rostea încă aceste cuvinte, Duhul Sfânt a căzut peste toți cei ce auzeau cuvântul.

\VS{45}Cei din circumcizie care au crezut erau uimiți, toți cei care veniseră cu Petru, pentru că darul Duhului Sfânt se revărsase și peste neamuri.

\VS{46}Căci îi auzeau vorbind în alte limbi și lăudând pe Dumnezeu.

\par }{\PP Atunci Petru a răspuns:

\VS{47}“Poate cineva să interzică acestor oameni să fie botezați cu apă? Ei au primit Duhul Sfânt la fel ca și noi”.

\VS{48}Și le-a poruncit să fie botezați în numele lui Isus Hristos. Apoi l-au rugat să rămână câteva zile.



\par }\Chap{11}{\PP \VerseOne{1}Apostolii și frații care erau în Iudeea au auzit că și neamurile au primit cuvântul lui Dumnezeu.

\VS{2}După ce Petru s-a suit la Ierusalim, cei care erau din circumcizie s-au certat cu el,

\VS{3}spunând: “Ai intrat la oameni netăiați împrejur și ai mâncat cu ei!”.


\par }{\PP \VS{4}Dar Petru a început și le-a explicat în ordine, zicând:

\VS{5}“Eram în cetatea Iope, rugându-mă, și am avut o vedenie: un vas care se cobora ca un cearșaf mare, care se lăsa din cer prin patru colțuri. A ajuns până la mine.

\VS{6}După ce m-am uitat cu atenție la el, m-am gândit și am văzut animalele cu patru picioare de pe pământ, animalele sălbatice, târâtoarele și păsările cerului.

\VS{7}Am auzit de asemenea un glas care îmi spunea:
{\WJ{“Scoală-te, Petru, ucide și mănâncă!”
}}

\VS{8}Dar eu am zis: “Nu așa, Doamne, căci în gura mea nu a intrat niciodată nimic nesfânt sau necurat.

\VS{9}Dar un glas mi-a răspuns a doua oară din cer: “{\WJ{Ceea ce Dumnezeu a curățat, tu să nu numești necurat”.
}}

\VS{10}Aceasta s-a făcut de trei ori și toți au fost trași din nou în cer.

\VS{11}Iată că imediat au stat în fața casei în care mă aflam trei bărbați, trimiși de la Cezareea la mine.

\VS{12}Duhul mi-a spus să merg cu ei fără să fac deosebire. M-au însoțit și acești șase frați și am intrat în casa acelui om.

\VS{13}El ne-a povestit cum îl văzuse pe înger stând în casa lui și spunându-i: “Trimite la Iope și adu-l pe Simon, care se numește Petru,

\VS{14}care îți va spune cuvintele prin care vei fi mântuit, tu și toată casa ta.

\VS{15}Când am început să vorbesc, Duhul Sfânt a căzut peste ei, ca și peste noi la început.

\VS{16}Mi-am adus aminte de cuvântul Domnului, cum a spus:
{\WJ{'Ioan a botezat într-adevăr în apă, dar voi veți fi botezați în Duhul Sfânt'.
}}

\VS{17}Dacă deci Dumnezeu le-a dat același dar ca și nouă, când am crezut în Domnul Isus Hristos, cine eram eu, ca să mă împotrivesc lui Dumnezeu?”


\par }{\PP \VS{18}Când au auzit acestea, au tăcut și au slăvit pe Dumnezeu, zicând: “Dumnezeu a dat și neamurilor pocăința pentru viață!”


\par }{\PP \VS{19}De aceea, cei ce se împrăștiaseră din pricina asupririi lui Ștefan au călătorit până în Fenicia, în Cipru și în Antiohia, și nu au vorbit nimănui decât numai iudeilor.

\VS{20}Dar au fost unii dintre ei, bărbați din Cipru și din Cirene, care, după ce au ajuns în Antiohia, au vorbit elinilor,
\FTNT{*}{{\FR 11:20 }{\FT{Un elenist este o persoană care păstrează obiceiurile și cultura greacă.}}}propovăduind pe Domnul Isus.

\VS{21}Mâna Domnului a fost cu ei și un mare număr a crezut și s-a întors la Domnul.

\VS{22}Vestea despre ei a ajuns la urechile adunării care era la Ierusalim. Ei l-au trimis pe Barnaba să meargă până la Antiohia,

\VS{23}care, după ce a ajuns și a văzut harul lui Dumnezeu, s-a bucurat. El i-a îndemnat pe toți, ca, cu hotărâre de inimă, să rămână aproape de Domnul.

\VS{24}Căci era un om bun, plin de Duhul Sfânt și de credință, și mulți oameni s-au adăugat la Domnul.


\par }{\PP \VS{25}Barnaba a plecat la Tars, ca să caute pe Saul.

\VS{26}După ce l-a găsit, l-a adus la Antiohia. Timp de un an întreg au fost adunați împreună cu adunarea și au învățat mult popor. Discipolii au fost numiți pentru prima dată creștini în Antiohia.


\par }{\PP \VS{27}În zilele acestea, au coborât prooroci din Ierusalim în Antiohia.

\VS{28}Unul dintre ei, pe nume Agabus, s-a ridicat în picioare și a arătat prin Duhul Sfânt că va fi o mare foamete în toată lumea, ceea ce s-a întâmplat și în zilele lui Claudius.

\VS{29}Cum vreunul dintre discipoli avea belșug, fiecare a hotărât să trimită ajutoare fraților care locuiau în Iudeea;

\VS{30}ceea ce au și făcut, trimițându-le bătrânilor prin mâna lui Barnaba și a lui Saul.



\par }\Chap{12}{\PP \VerseOne{1}În vremea aceea, regele Irod și-a întins mâinile ca să asuprească pe unii din adunare.

\VS{2}Și a ucis cu sabia pe Iacov, fratele lui Ioan.

\VS{3}Când a văzut că aceasta a fost pe placul iudeilor, s-a apucat să prindă și pe Petru. Acest lucru se întâmpla în zilele azimilor.

\VS{4}După ce l-a arestat, l-a băgat în închisoare și l-a predat la patru escadroane de câte patru soldați fiecare, ca să-l păzească, cu intenția de a-l scoate în fața poporului după Paști.

\VS{5}Petru era deci ținut în închisoare, dar adunarea se ruga neîncetat lui Dumnezeu pentru el.

\VS{6}În noaptea în care Irod era pe cale să-l scoată afară, Petru dormea între doi soldați, legat cu două lanțuri. Gărzile din fața ușii păzeau închisoarea.


\par }{\PP \VS{7}Și iată că un înger al Domnului stătea lângă el, și o lumină strălucea în chilie. L-a lovit pe Petru peste o parte și l-a trezit, spunându-i: “Ridică-te repede!”. Lanțurile i-au căzut de pe mâini.

\VS{8}Îngerul i-a spus: “Îmbracă-te și pune-ți sandalele”. Așa a făcut el. El i-a spus: “Pune-ți mantia și urmează-mă”.

\VS{9}Și el a ieșit și l-a urmat. El nu știa că ceea ce făcea îngerul era real, ci credea că a văzut o viziune.

\VS{10}După ce au trecut de prima și de a doua gardă, au ajuns la poarta de fier care duce în cetate, care li s-a deschis singură. Au ieșit și au coborât pe o stradă, și imediat îngerul s-a îndepărtat de el.


\par }{\PP \VS{11}După ce și-a revenit, Petru a zis: “Acum știu cu adevărat că Domnul a trimis pe îngerul Său și m-a izbăvit din mâna lui Irod și de tot ce aștepta poporul iudeu.”

\VS{12}Gândindu-se la aceasta, a venit la casa Mariei, mama lui Ioan, zis Marcu, unde erau mulți adunați și se rugau.

\VS{13}Când Petru a bătut la ușa porții, o slujnică pe nume Rodica a venit să răspundă.

\VS{14}Când a recunoscut vocea lui Petru, nu a deschis poarta de bucurie, ci a alergat înăuntru și a anunțat că Petru stătea în fața porții.


\par }{\PP \VS{15}Ei i-au zis: “Ești nebună!” Dar ea insista că așa este. Ei au spus: “Este îngerul lui.”

\VS{16}Dar Petru a continuat să bată la ușă. După ce au deschis, l-au văzut și au rămas uimiți.

\VS{17}Dar el, făcându-le semn cu mâna să tacă, le-a povestit cum l-a scos Domnul din închisoare. Și a zis: “Spuneți aceste lucruri lui Iacov și fraților”. Apoi a plecat și s-a dus în alt loc.


\par }{\PP \VS{18}Și cum s-a făcut ziuă, nu mică era agitația printre ostași despre ce se întâmplase cu Petru.

\VS{19}După ce Irod l-a căutat și nu l-a găsit, a cercetat gărzile, apoi a poruncit să fie omorât. S-a coborât din Iudeea în Cezareea și a rămas acolo.


\par }{\PP \VS{20}Irod era foarte mânios pe poporul din Tir și din Sidon. Aceștia au venit de comun acord la el și, făcându-l prieten pe Blastus, ajutorul personal al regelui, i-au cerut pace, pentru că țara lor depindea de țara regelui pentru hrană.

\VS{21}În ziua stabilită, Irod s-a îmbrăcat în haine regale, s-a așezat pe tron și le-a ținut un discurs.

\VS{22}Poporul a strigat: “Vocea unui zeu, și nu a unui om!”

\VS{23}Imediat, un înger al Domnului l-a lovit, pentru că nu i-a dat slavă lui Dumnezeu. Apoi a fost mâncat de viermi și a murit.


\par }{\PP \VS{24}Dar cuvântul lui Dumnezeu creștea și se înmulțea.

\VS{25}Barnaba și Saul s-au întors la
\FTNT{*}{{\FR 12:25 }{\FT{TR citește “de la” în loc de “la”}}}Ierusalim, după ce și-au împlinit slujba, luând cu ei și pe Ioan, numit Marcu.



\par }\Chap{13}{\PP \VerseOne{1}În adunarea care era în Antiohia erau niște prooroci și învățători: Barnaba, Simeon, care se numea Niger, Lucius din Cirene, Manaen, fratele adoptiv al tetrarhului Irod, și Saul.

\VS{2}În timp ce slujeau Domnului și posteau, Duhul Sfânt a zis: “Separă-i pe Barnaba și pe Saul pentru mine, pentru lucrarea la care i-am chemat”.


\par }{\PP \VS{3}Apoi, după ce au postit și s-au rugat și și-au pus mâinile peste ei, i-au trimis.

\VS{4}Astfel, fiind trimiși de Duhul Sfânt, s-au coborât la Seleucia. De acolo au navigat spre Cipru.

\VS{5}Când au ajuns la Salamina, au vestit cuvântul lui Dumnezeu în sinagogile evreiești. L-au avut și pe Ioan ca însoțitor.

\VS{6}După ce au străbătut insula până la Pafos, au găsit un vrăjitor, un fals profet, un iudeu care se numea Bar Isus,

\VS{7}care se afla cu proconsulul Sergius Paulus, un om înțelegător. Acesta i-a chemat pe Barnaba și pe Saul și a căutat să audă cuvântul lui Dumnezeu.

\VS{8}Dar vrăjitorul Elymas (căci așa se numește el după interpretare) s-a împotrivit lor, căutând să-l întoarcă pe proconsul de la credință.

\VS{9}Dar Saul, care se mai numește și Pavel, plin de Duhul Sfânt, și-a fixat ochii asupra lui

\VS{10}și a zis: “Fiu al diavolului, plin de toată înșelăciunea și de toată viclenia, vrăjmaș al oricărei dreptăți, nu vei înceta să pervertești căile drepte ale Domnului?

\VS{11}Acum, iată, mâna Domnului este asupra ta și vei fi orb, fără să vezi soarele pentru o vreme!”

\par }{\PP Imediat, o ceață și întunericul au căzut peste el. A umblat de colo-colo căutând pe cineva care să-l ducă de mână.

\VS{12}Atunci proconsulul, când a văzut ce se întâmplase, a crezut, fiind uimit de învățătura Domnului.


\par }{\PP \VS{13}Pavel și ceata lui au pornit din Pafos și au ajuns la Perga, în Pamfilia. Ioan s-a despărțit de ei și s-a întors la Ierusalim.

\VS{14}Dar ei, trecând de la Perga, au ajuns la Antiohia Pisidiei. Au intrat în sinagogă în ziua de Sabat și au șezut jos.

\VS{15}După citirea Legii și a profeților, conducătorii sinagogii au trimis la ei, spunându-le: “Fraților, dacă aveți vreun cuvânt de îndemn pentru popor, vorbiți.”


\par }{\PP \VS{16}Pavel s-a ridicat în picioare și, făcând semn cu mâna, a zis: “Bărbați israeliți și voi, cei ce vă temeți de Dumnezeu, ascultați.

\VS{17}Dumnezeul acestui popor i-a
\FTNT{*}{{\FR 13:17 }{\FT{TR, NU adaugă “Israel”}}}ales pe părinții noștri și a înălțat poporul când a rămas ca străin în țara Egiptului și, cu brațul ridicat, l-a scos din ea.

\VS{18}Pe o perioadă de aproximativ patruzeci de ani i-a suportat în pustiu.

\VS{19}După ce a distrus șapte națiuni în țara Canaanului, le-a dat țara lor ca moștenire timp de aproximativ patru sute cincizeci de ani.

\VS{20}După aceste lucruri, le-a dat judecători până la profetul Samuel.

\VS{21}După aceea au cerut un rege, și Dumnezeu le-a dat pe Saul, fiul lui Chiș, un bărbat din tribul lui Beniamin, timp de patruzeci de ani.

\VS{22}După ce l-a înlăturat, l-a ridicat ca rege al lor pe David, căruia i-a și mărturisit: “Am găsit pe David, fiul lui Isai, un om după inima mea, care va face toată voia mea.

\VS{23}Din urmașul acestui om, Dumnezeu a adus mântuirea\FTNT{†}{{\FR 13:23 }{\FT{TR, NU citiți “un Mântuitor, Isus” în loc de “mântuire”}}} lui Israel, conform promisiunii sale,

\VS{24}înainte de venirea sa, când Ioan a predicat pentru prima dată lui Israel botezul pocăinței.
\FTNT{‡}{{\FR 13:24 }{\FT{TR, NU citiți “la tot poporul lui Israel” în loc de “la Israel”}}}

\VS{25}Pe când Ioan își împlinea cursul, a zis: “Ce credeți că sunt eu? Eu nu sunt el. Dar iată că vine după mine unul, ale cărui sandale nu sunt vrednic să le dezleg'.


\par }{\PP \VS{26}“Fraților, copii din neamul lui Avraam, și cei din voi care vă temeți de Dumnezeu, cuvântul acestei mântuiri a fost trimis la voi.

\VS{27}Căci cei care locuiesc în Ierusalim și conducătorii lor, pentru că nu l-au cunoscut, nici glasurile profeților care se citesc în fiecare Sabat, le-au împlinit condamnându-l.

\VS{28}Deși nu au găsit nici un motiv de moarte, i-au cerut totuși lui Pilat să-L omoare.

\VS{29}După ce au împlinit tot ce era scris despre el, l-au dat jos de pe lemn și l-au pus într-un mormânt.

\VS{30}Dar Dumnezeu l-a înviat din morți,

\VS{31}și a fost văzut timp de multe zile de cei care au urcat cu el din Galileea la Ierusalim, care sunt martorii lui în fața poporului.

\VS{32}Noi vă aducem vestea cea bună a promisiunii făcute părinților,

\VS{33}că Dumnezeu a împlinit-o pentru noi, copiii lor, prin faptul că l-a înviat pe Isus. După cum este scris și în al doilea psalm,

\par }{\Q “Tu ești Fiul meu.

\par }{\QB Astăzi am devenit tatăl tău”.
\XREF{✡}{{\XO 13:33 }{\XT{{\CHAR{Psalmul 2:7}}}}}


\par }{\PP \VS{34}“Despre faptul că l-a înviat din morți, ca să nu se mai întoarcă la stricăciune, așa a vorbit: 'Vă voi da binecuvântările sfinte și sigure ale lui David'.
\XREF{✡}{{\XO 13:34 }{\XT{{\CHAR{Isaia 55:3}}}}}

\VS{35}De aceea spune și într-un alt psalm: “Nu vei îngădui ca Sfântul Tău să vadă putrezirea”.
\XREF{✡}{{\XO 13:35 }{\XT{{\CHAR{Psalmul 16:10}}}}}

\VS{36}Căci David, după ce a slujit în neamul său sfatul lui Dumnezeu, a adormit, a fost culcat cu părinții săi și a văzut putrezirea.

\VS{37}Dar cel pe care l-a înviat Dumnezeu nu a văzut putrezirea.

\VS{38}Să știți, deci, fraților, că prin acest om vi se vestește iertarea păcatelor;

\VS{39}și prin el oricine crede este îndreptățit de toate lucrurile de care nu ați putut fi îndreptățiți prin Legea lui Moise.

\VS{40}Feriți-vă, deci, ca să nu vină peste voi ceea ce s-a spus în profeți:


\par }{\Q \VS{41}“Iată, batjocoritorilor!

\par }{\QB Minunează-te și pieri,

\par }{\QB Căci Eu fac o lucrare în zilele voastre,

\par }{\QB o lucrare pe care nu o veți crede în nici un fel, dacă cineva v-o declară.””
\FTNT{§}{{\FR 13:41 }{\FT{Habacuc 1:5}}}


\par }{\PP \VS{42}Iudeii au ieșit din sinagogă, și neamurile au cerut să li se vestească aceste cuvinte în Sabatul următor.

\VS{43}După ce s-a despărțit sinagoga, mulți dintre iudei și dintre prozeliții evlavioși au urmat pe Pavel și pe Barnaba, care, vorbindu-le, îi îndemnau să continue în harul lui Dumnezeu.


\par }{\PP \VS{44}În sâmbăta următoare, aproape toată cetatea s-a adunat să asculte cuvântul lui Dumnezeu.

\VS{45}Dar iudeii, văzând mulțimea, s-au umplut de gelozie, au contrazis cele spuse de Pavel și au hulit.


\par }{\PP \VS{46}Pavel și Barnaba au luat cuvântul cu îndrăzneală și au zis: “Trebuia să vă fie vestit mai întâi cuvântul lui Dumnezeu. Întrucât, într-adevăr, îl respingeți de la voi înșivă și vă judecați nevrednici de viața veșnică, iată că ne întoarcem la neamuri.

\VS{47}Căci așa ne-a poruncit Domnul, zicând:,

\par }{\Q “Te-am pus ca lumină pentru neamuri,

\par }{\QB ca să aduci mântuirea până la marginile pământului.”
\FTNT{*}{{\FR 13:47 }{\FT{{\CHAR{Isaia 49:6}}}}}


\par }{\PP \VS{48}Neamurile, auzind acestea, s-au bucurat și au slăvit cuvântul lui Dumnezeu. Toți cei care erau rânduiți pentru viața veșnică au crezut.

\VS{49}Cuvântul Domnului s-a răspândit în toată regiunea.

\VS{50}Dar iudeii au stârnit pe femeile evlavioase și de seamă și pe căpeteniile cetății și au stârnit o persecuție împotriva lui Pavel și a lui Barnaba și i-au aruncat afară din hotarele lor.

\VS{51}Dar ei și-au scuturat praful de pe picioare împotriva lor și au venit la Iconia.

\VS{52}Ucenicii au fost plini de bucurie și de Duhul Sfânt.



\par }\Chap{14}{\PP \VerseOne{1}În Iconiu, au intrat împreună în sinagoga iudeilor și au vorbit astfel încât o mare mulțime de iudei și de greci au crezut.

\VS{2}Dar iudeii necredincioși\FTNT{*}{{\FR 14:2 }{\FT{sau, neascultător}}} au stârnit și au învrăjbit sufletele neamurilor împotriva fraților.

\VS{3}De aceea au rămas acolo multă vreme, vorbind cu îndrăzneală în Domnul, care dădea mărturie despre cuvântul harului său, îngăduind să se facă semne și minuni prin mâinile lor.

\VS{4}Dar mulțimea din cetate era împărțită. O parte era de partea iudeilor și o parte de partea apostolilor.

\VS{5}Când o parte dintre neamuri și iudei, împreună cu conducătorii lor, au încercat cu violență să-i maltrateze și să-i ucidă cu pietre,

\VS{6}ei au luat cunoștință de aceasta și au fugit în orașele Licaonia, Listra, Derbe și în regiunea înconjurătoare.

\VS{7}Acolo au propovăduit Vestea cea Bună.


\par }{\PP \VS{8}La Listra ședea un om neputincios la picioare, olog din pântecele mamei sale, care nu umblase niciodată.

\VS{9}El asculta vorbind pe Pavel, care, fixând ochii asupra lui și văzând că are credința de a fi vindecat,

\VS{10}a zis cu glas tare: “Ridică-te în picioare!” El a sărit în picioare și a mers.

\VS{11}Când mulțimea a văzut ce făcuse Pavel, a ridicat glasul și a zis în limba Licaoniei: “Zeii s-au coborât la noi în chip de oameni!”

\VS{12}Ei îl numeau pe Barnaba “Jupiter”, iar pe Pavel “Mercur”, pentru că el era cel care vorbea cel mai mult.

\VS{13}Preotul lui Jupiter, al cărui templu se afla în fața cetății lor, a adus boi și ghirlande la porți și ar fi vrut să facă un sacrificiu împreună cu mulțimea.


\par }{\PP \VS{14}Apostolii Barnaba și Pavel, auzind, și-au rupt hainele și au sărit în mulțime, strigând:

\VS{15}“Bărbaților, pentru ce faceți aceste lucruri? Și noi suntem oameni de aceeași natură ca și voi și vă aducem o veste bună, ca să vă întoarceți de la aceste lucruri deșarte la Dumnezeul cel viu, care a făcut cerul, pământul, marea și tot ce este în ele;

\VS{16}care, în generațiile trecute, a lăsat toate neamurile să umble pe căile lor.

\VS{17}Și totuși nu s-a lăsat fără mărturie, căci a făcut binele și v-a\FTNT{†}{{\FR 14:17 }{\FT{TR citește “noi” în loc de “voi”}}} dat ploi din cer și anotimpuri roditoare, umplându-ne inimile de hrană și de bucurie.”


\par }{\PP \VS{18}Și chiar dacă spuneau aceste lucruri, cu greu au împiedicat mulțimile să le aducă jertfe.

\VS{19}Dar au venit acolo niște iudei din Antiohia și din Iconiu și, după ce au convins mulțimile, l-au ucis cu pietre pe Pavel și l-au târât afară din cetate, crezând că este mort.


\par }{\PP \VS{20}Dar, pe când stăteau ucenicii în jurul Lui, S-a sculat și a intrat în cetate. A doua zi, a ieșit cu Barnaba la Derbe.


\par }{\PP \VS{21}După ce au propovăduit vestea cea bună în cetatea aceea și au făcut mulți ucenici, s-au întors la Listra, la Iconiu și la Antiohia,

\VS{22}și au întărit sufletele ucenicilor, îndemnându-i să rămână în credință și să le spună că prin multe necazuri trebuie să intrăm în Împărăția lui Dumnezeu.

\VS{23}După ce au numit pentru ei prezbiteri în fiecare adunare și s-au rugat cu post, i-au încredințat Domnului în care crezuseră.


\par }{\PP \VS{24}Au trecut prin Pisidia și au ajuns în Pamfilia.

\VS{25}După ce au rostit cuvântul în Perga, s-au coborât în Attalia.

\VS{26}De acolo au navigat spre Antiohia, de unde fuseseră încredințați harului lui Dumnezeu pentru lucrarea pe care o îndepliniseră.

\VS{27}După ce au ajuns și au adunat adunarea, au relatat toate lucrurile pe care Dumnezeu le făcuse cu ei și faptul că deschisese o ușă a credinței pentru națiuni.

\VS{28}Au rămas acolo cu discipolii pentru o lungă perioadă de timp.



\par }\Chap{15}{\PP \VerseOne{1}Niște oameni coborâți din Iudeea au învățat pe frați “Dacă nu vă tăiați împrejur după obiceiul lui Moise, nu puteți fi mântuiți.”

\VS{2}De aceea, după ce Pavel și Barnaba au avut cu ei nu puține neînțelegeri și discuții, au însărcinat pe Pavel, pe Barnaba și pe alții dintre ei să se suie la Ierusalim la apostoli și la bătrâni în legătură cu această chestiune.

\VS{3}Ei, fiind trimiși pe drum de adunare, au trecut prin Fenicia și Samaria, vestind convertirea neamurilor. Ei au provocat o mare bucurie tuturor fraților.

\VS{4}După ce au ajuns la Ierusalim, au fost primiți de Adunare, de apostoli și de bătrâni și au relatat tot ce făcuse Dumnezeu cu ei.


\par }{\PP \VS{5}Dar unii din secta fariseilor credincioși s-au sculat și au zis: “Trebuie să-i tăiem împrejur și să le poruncim să păzească Legea lui Moise.”


\par }{\PP \VS{6}Și apostolii și bătrânii s-au adunat ca să vadă despre aceasta.

\VS{7}După ce s-a discutat mult, Petru s-a ridicat și le-a zis: “Fraților, știți că, cu mult timp în urmă, Dumnezeu a făcut o alegere între voi, pentru ca, prin gura mea, neamurile să audă cuvântul Bunei Vestiri și să creadă.

\VS{8}Dumnezeu, care cunoaște inima, a mărturisit despre ei, dându-le Duhul Sfânt, așa cum ne-a făcut și nouă.

\VS{9}El nu a făcut nicio deosebire între noi și ei, curățându-le inimile prin credință.

\VS{10}Așadar, acum, de ce îl ispitiți pe Dumnezeu, ca să puneți pe gâtul discipolilor un jug pe care nici părinții noștri și nici noi nu am putut să-l purtăm?

\VS{11}Noi, însă,\FTNT{*}{{\FR 15:11 }{\FT{TR adaugă “Hristos”}}} credem că suntem mântuiți prin harul Domnului Isus, ca și ei.”


\par }{\PP \VS{12}Toată mulțimea tăcea și asculta pe Barnaba și pe Pavel, care vorbeau despre semnele și minunile pe care le făcuse Dumnezeu printre neamuri prin ei.

\VS{13}După ce au tăcut, Iacov a răspuns: “Fraților, ascultați-mă.

\VS{14}Simeon a relatat cum Dumnezeu a vizitat mai întâi națiunile pentru a lua din ele un popor pentru numele său.

\VS{15}Acest lucru este în acord cu cuvintele profeților. După cum este scris,


\par }{\Q \VS{16}“După acestea mă voi întoarce.

\par }{\Q Voi zidi din nou cortul lui David, care a căzut.

\par }{\Q Voi construi din nou ruinele sale.

\par }{\Q Îl voi ridica

\VS{17}pentru ca ceilalți oameni să caute pe Domnul:

\par }{\Q toate neamurile care sunt chemate după numele Meu,

\par }{\Q zice Domnul, care face toate aceste lucruri.
\FTNT{†}{{\FR 15:17 }{\FT{Amos 9:11-12}}}


\par }{\PP \VS{18}“Toate lucrările lui Dumnezeu sunt cunoscute de El din veșnicie.

\VS{19}De aceea, judecata mea este să nu-i tulburăm pe cei din rândul neamurilor care se întorc la Dumnezeu,

\VS{20}ci să le scriem să se abțină de la spurcarea idolilor, de la imoralitatea sexuală, de la ceea ce este sugrumat și de la sânge.

\VS{21}Căci Moise, din generații de demult, are în fiecare cetate pe cei care îl propovăduiesc, fiind citit în sinagogi în fiecare Sabat.”


\par }{\PP \VS{22}Atunci apostolii și bătrânii, împreună cu toată adunarea, au găsit de cuviință să aleagă bărbați din ceata lor și să-i trimită la Antiohia cu Pavel și Barnaba: Iuda, numit Barsaba, și Sila, bărbați de seamă între frați.

\VS{23}Ei au scris aceste lucruri de mâna lor:

\par }{\PP “Apostolii, bătrânii și frații, către frații care sunt dintre neamuri în Antiohia, Siria și Cilicia: salutări.

\VS{24}Pentru că am auzit că unii care au plecat de la noi v-au tulburat cu vorbe, neliniștindu-vă sufletele, spunând: “Trebuie să vă tăiați împrejur și să țineți Legea”, cărora noi nu le-am dat nicio poruncă;

\VS{25}ni s-a părut bine, după ce am ajuns la un acord, să alegem niște oameni și să vi-i trimitem împreună cu preaiubiții noștri Barnaba și Pavel,

\VS{26}oameni care și-au riscat viața pentru numele Domnului nostru Isus Hristos.

\VS{27}Am trimis, așadar, pe Iuda și pe Sila, care și ei înșiși vă vor spune aceleași lucruri prin cuvânt.

\VS{28}Fiindcă Duhul Sfânt și nouă ni s-a părut bine să nu punem asupra voastră o povară mai mare decât aceste lucruri necesare:

\VS{29}să vă abțineți de la lucrurile jertfite idolilor, de la sânge, de la lucrurile sugrumate și de la imoralitatea sexuală, de care, dacă vă veți păzi, vă va fi bine. Adio!”


\par }{\PP \VS{30}După ce au fost trimiși, au ajuns la Antiohia. După ce au adunat mulțimea, au înmânat scrisoarea.

\VS{31}După ce au citit-o, s-au bucurat de încurajare.

\VS{32}Iuda și Sila, fiind și ei prooroci, i-au încurajat pe frați cu multe cuvinte și i-au întărit.

\VS{33}După ce au petrecut ceva timp acolo, frații i-au concediat în pace de la frați la apostoli.

\VS{34}\FTNT{‡}{{\FR 15:34 }{\FT{Unele manuscrise adaugă:
}}{\FQA{Dar lui Sila i s-a părut bine să rămână acolo.
}}}

\VS{35}Dar Pavel și Barnaba au rămas în Antiohia, învățând și propovăduind cuvântul Domnului, împreună cu mulți alții.


\par }{\PP \VS{36}După câteva zile, Pavel a zis lui Barnaba: “Să ne întoarcem acum și să vizităm pe frații noștri în fiecare cetate în care am vestit cuvântul Domnului, ca să vedem ce mai fac.”

\VS{37}Barnaba plănuia să ia cu ei și pe Ioan, căruia i se spunea Marcu.

\VS{38}Dar lui Pavel nu i s-a părut o idee bună să ia cu ei pe cineva care se retrăsese de la ei în Pamfilia și care nu a mers cu ei să facă lucrarea.

\VS{39}Atunci cearta a devenit atât de aprinsă, încât s-au despărțit unul de altul. Barnaba l-a luat pe Marcu cu el și a plecat în Cipru,

\VS{40}dar Pavel l-a ales pe Sila și a plecat, fiind încredințat de frați harului lui Dumnezeu.

\VS{41}A străbătut Siria și Cilicia, întărind adunările.



\par }\Chap{16}{\PP \VerseOne{1}A venit la Derbe și la Listra; și iată că era acolo un ucenic, numit Timotei, fiul unei iudeice credincioase, dar tatăl său era grec.

\VS{2}Frații care erau la Listra și la Iconiu au dat o mărturie bună despre el.

\VS{3}Pavel a vrut ca el să meargă cu el; l-a luat și l-a circumcis din cauza iudeilor care erau în acele părți, căci toți știau că tatăl lui era grec.

\VS{4}Pe când mergeau prin cetăți, le transmiteau decretele pe care trebuiau să le respecte și care fuseseră rânduite de apostolii și de bătrânii care erau la Ierusalim.

\VS{5}Astfel, adunările se întăreau în credință și se înmulțeau în fiecare zi.


\par }{\PP \VS{6}După ce au străbătut ținutul Frigiei și al Galatiei, Duhul Sfânt le-a interzis să vorbească în Asia.

\VS{7}După ce au ajuns vizavi de Mysia, au încercat să intre în Bitinia, dar Duhul Sfânt nu le-a permis.

\VS{8}Trecând de Mysia, au coborât la Troa.

\VS{9}O viziune i-a apărut lui Pavel în timpul nopții. Era un om din Macedonia care stătea în picioare, îl implora și îi spunea: “Treci în Macedonia și ajută-ne”.

\VS{10}După ce a văzut viziunea, imediat am căutat să mergem în Macedonia, ajungând la concluzia că Domnul ne-a chemat să le vestim Vestea cea Bună.

\VS{11}Așadar, pornind din Troa, am pornit din Troa, am făcut drum drept spre Samotracia, iar a doua zi spre Neapolis;

\VS{12}și de acolo spre Filipi, care este un oraș din Macedonia, cel mai important din district, o colonie romană. Am stat câteva zile în acest oraș.


\par }{\PP \VS{13}În ziua de Sabat, ne-am dus în afara cetății, la marginea unui râu, unde am crezut că este un loc de rugăciune, și am luat loc și am vorbit cu femeile care se adunaseră.

\VS{14}Ne-a ascultat o femeie numită Lidia, vânzătoare de purpură, din cetatea Tiatira, o femeie care se închina lui Dumnezeu. Domnul i-a deschis inima ca să asculte cele spuse de Pavel.

\VS{15}După ce ea și familia ei au fost botezate, ne-a rugat, zicând: “Dacă ați judecat că sunt credincioasă Domnului, intrați în casa mea și rămâneți.” Așa ne-a convins ea.


\par }{\PP \VS{16}Pe când ne duceam la rugăciune, ne-a întâlnit o fată cu duh de ghicitoare, care aducea stăpânilor ei mult câștig prin ghicire.

\VS{17}Urmărindu-ne pe Pavel și pe noi, ea a strigat: “Acești oameni sunt slujitori ai Dumnezeului Celui Preaînalt, care ne vestesc o cale de mântuire!”

\VS{18}Ea a făcut acest lucru timp de mai multe zile.

\par }{\PP Dar Pavel, foarte supărat, s-a întors și i-a spus duhului: “Îți poruncesc în numele lui Isus Hristos să ieși din ea!”. A ieșit chiar în acea oră.

\VS{19}Dar când stăpânii ei au văzut că speranța câștigului lor a dispărut, i-au prins pe Pavel și pe Sila și i-au târât în piață în fața conducătorilor.

\VS{20}După ce i-au adus în fața magistraților, aceștia au spus: “Acești oameni, fiind iudei, ne agită cetatea

\VS{21}și susțin obiceiuri pe care nouă, fiind romani, nu ne este îngăduit să le acceptăm sau să le respectăm.”


\par }{\PP \VS{22}Mulțimea s-a ridicat împotriva lor, iar magistrații le-au rupt hainele și au poruncit să fie bătuți cu toiege.

\VS{23}După ce le-au aplicat multe lovituri, i-au aruncat în închisoare, însărcinând pe temnicer să-i păzească în siguranță.

\VS{24}După ce a primit o astfel de poruncă, i-a aruncat în închisoarea interioară și le-a asigurat picioarele în lanțuri.


\par }{\PP \VS{25}Dar, pe la miezul nopții, Pavel și Sila se rugau și cântau cântări lui Dumnezeu, iar cei închiși îi ascultau.

\VS{26}Deodată s-a produs un cutremur mare, încât s-au zguduit temeliile închisorii; și îndată s-au deschis toate ușile, și toate legăturile tuturor au fost desfăcute.

\VS{27}Temnicerul, fiind trezit din somn și văzând ușile închisorii deschise, și-a scos sabia și era pe punctul de a se sinucide, crezând că prizonierii au scăpat.

\VS{28}Dar Pavel a strigat cu glas tare: “Nu-ți face rău, căci suntem cu toții aici!”


\par }{\PP \VS{29}A chemat lumină, a intrat, a sărit înăuntru, a căzut tremurând înaintea lui Pavel și a lui Sila,

\VS{30}i-a scos afară și a zis: “Domnilor, ce trebuie să fac ca să mă mântuiesc?”


\par }{\PP \VS{31}Ei au zis: “Credeți în Domnul Isus Hristos și veți fi mântuiți, tu și casa ta.”

\VS{32}Ei i-au spus cuvântul Domnului lui și tuturor celor care erau în casa lui.


\par }{\PP \VS{33}La aceeași oră din noapte, i-a luat și le-a spălat rănile, și îndată s-a botezat el și toată casa lui.

\VS{34}I-a urcat în casa lui, le-a pus mâncare înaintea lor și s-a bucurat foarte mult împreună cu toată casa lui, pentru că a crezut în Dumnezeu.


\par }{\PP \VS{35}Dar, când s-a făcut ziuă, magistrații au trimis pe sergenți, zicând: “Lasă-i să plece pe oamenii aceia.”


\par }{\PP \VS{36}Temnicerul a raportat aceste cuvinte lui Pavel, zicând: “Magistrații au trimis să te lase să pleci; ieși acum și du-te în pace.”


\par }{\PP \VS{37}Dar Pavel le-a zis: “Ne-au bătut în public, fără judecată, niște romani, și ne-au aruncat în temniță! Ne eliberează ei acum în secret? Nu, cu siguranță, ci să vină ei înșiși și să ne scoată afară!”


\par }{\PP \VS{38}Și sergenții au raportat aceste cuvinte magistraților, care s-au temut când au auzit că erau romani,

\VS{39}și au venit și i-au rugat. După ce i-au scos afară, i-au rugat să plece din oraș.

\VS{40}Au ieșit din închisoare și au intrat în casa Lidiei. După ce i-au văzut pe frați, i-au încurajat, apoi au plecat.



\par }\Chap{17}{\PP \VerseOne{1}După ce au trecut prin Amfipoli și Apolonia, au ajuns la Tesalonic, unde era o sinagogă iudaică.

\VS{2}Pavel, după obiceiul său, a intrat la ei și, timp de trei zile de Sabat, a discutat cu ei din Scripturi,

\VS{3}explicând și demonstrând că Hristos trebuia să sufere și să învieze din morți și spunând: “Acest Isus, pe care vi-L vestesc, este Hristosul”.


\par }{\PP \VS{4}Unii dintre ei au fost convinși și s-au alăturat lui Pavel și Sila; dintre greci, o mare mulțime de credincioși și nu puține femei de seamă.

\VS{5}Dar iudeii neînduplecați au luat cu\FTNT{*}{{\FR 17:5 }{\FT{TR citește “Și iudeii care nu erau convinși, devenind invidioși și luându-i cu ei” în loc de “Dar iudeii care nu erau convinși i-au luat cu ei”.}}} ei câțiva oameni răi din piață și, adunând o mulțime, au pus cetatea în tumult. Asaltând casa lui Iason, au căutat să-i scoată în fața poporului.

\VS{6}Cum nu i-au găsit, i-au târât pe Iason și pe unii frați în fața conducătorilor cetății, strigând: “Au venit aici și aceștia care au răsturnat lumea,

\VS{7}pe care i-a primit Iason. Toți aceștia acționează împotriva hotărârilor Cezarului, spunând că există un alt rege, Isus!”.

\VS{8}Mulțimea și conducătorii cetății s-au tulburat când au auzit aceste lucruri.

\VS{9}După ce au luat garanții de la Iason și de la ceilalți, i-au lăsat să plece.


\par }{\PP \VS{10}Frații au trimis îndată pe Pavel și pe Sila noaptea la Beroe. Când au ajuns, au intrat în sinagoga iudaică.


\par }{\PP \VS{11}Aceștia erau mai nobili decât cei din Tesalonic, pentru că primeau Cuvântul cu toată disponibilitatea minții, cercetând zilnic Scripturile, ca să vadă dacă aceste lucruri sunt adevărate.

\VS{12}De aceea, mulți dintre ei au crezut; de asemenea, dintre femeile grecești de seamă și nu puțini bărbați.

\VS{13}Dar, când iudeii din Tesalonic au aflat că cuvântul lui Dumnezeu a fost vestit de Pavel și în Beroea, au venit și ei acolo, agitând mulțimile.

\VS{14}Atunci frații au trimis îndată pe Pavel să meargă până la mare, iar Sila și Timotei au rămas încă acolo.

\VS{15}Dar cei care îl însoțeau pe Pavel l-au dus până la Atena. Primind poruncă de la Sila și Timotei să vină foarte repede la el, au plecat.


\par }{\PP \VS{16}Pavel îi aștepta la Atena și, când a văzut cetatea plină de idoli, i s-a aprins duhul în el.

\VS{17}Așa că discuta în sinagogă cu iudeii și cu persoanele evlavioase și în piață, în fiecare zi, cu cei care îl întâlneau.

\VS{18}Unii dintre filosofii epicurieni și stoici stăteau de asemenea\FTNT{†}{{\FR 17:18 }{\FT{TR omite “de asemenea”
}}} de vorbă cu el. Unii ziceau: “Ce vrea să spună acest bălăcar?”.

\par }{\PP Alții au spus: “Se pare că susține zeități străine”, pentru că l-a predicat pe Isus și învierea.


\par }{\PP \VS{19}L-au prins și l-au dus în Areopag, zicând: “Să știm ce este această învățătură nouă, despre care vorbești?

\VS{20}Căci ne aduci la urechi niște lucruri ciudate. Vrem deci să știm ce înseamnă aceste lucruri”.

\VS{21}Or, toți atenienii și străinii care locuiau acolo nu-și petreceau timpul decât pentru a spune sau a auzi vreun lucru nou.


\par }{\PP \VS{22}Pavel, stând în mijlocul Areopagului, a zis: “Oameni din Atena, văd că sunteți foarte credincioși în toate.

\VS{23}Căci, trecând și observând obiectele cultului vostru, am găsit și un altar cu această inscripție: “PENTRU UN DUMNEZEU NECUNOSCUT.” De aceea, ceea ce voi adorați în necunoștință de cauză, vă anunț.

\VS{24}Dumnezeul care a făcut lumea și toate lucrurile din ea, el, care este Domnul cerului și al pământului, nu locuiește în temple făcute de mâini.

\VS{25}El nu este slujit de mâini omenești, ca și cum ar avea nevoie de ceva, întrucât el însuși dă tuturor viață, suflare și toate lucrurile.

\VS{26}El a făcut ca dintr-un singur sânge să locuiască orice neam de oameni pe toată suprafața pământului, stabilind anotimpurile și limitele locuințelor lor,

\VS{27}pentru ca ei să îl caute pe Domnul, dacă poate ar putea să întindă mâna spre el și să îl găsească, deși el nu este departe de fiecare dintre noi.

\VS{28}“Căci în El trăim, ne mișcăm și suntem”. Așa cum au spus unii dintre poeții voștri: 'Căci și noi suntem urmașii lui'.

\VS{29}Fiind, așadar, urmașii lui Dumnezeu, nu trebuie să credem că natura divină este ca aurul, argintul sau piatra, gravată prin arta și designul omului.

\VS{30}De aceea, Dumnezeu a trecut cu vederea vremurile de ignoranță. Dar acum poruncește ca toți oamenii de pretutindeni să se pocăiască,

\VS{31}pentru că a rânduit o zi în care va judeca lumea în dreptate, prin omul pe care l-a rânduit; despre care a dat asigurări tuturor oamenilor, prin faptul că l-a înviat din morți.”


\par }{\PP \VS{32}Iar când auzeau despre învierea morților, unii râdeau, iar alții ziceau: “Vrem să te mai auzim încă o dată despre aceasta.”


\par }{\PP \VS{33}Astfel a ieșit Pavel din mijlocul lor.

\VS{34}Dar unii oameni s-au alăturat lui și au crezut, printre care Dionisie Areopagitul, o femeie cu numele Damaris și alții împreună cu ei.



\par }\Chap{18}{\PP \VerseOne{1}După aceste lucruri, Pavel a plecat din Atena și a venit la Corint.

\VS{2}A găsit un iudeu numit Aquila, de neam din Pont, care venise de curând din Italia cu Priscila, soția sa, pentru că Claudius poruncise ca toți iudeii să plece din Roma. El a venit la ei

\VS{3}și, pentru că practica aceeași meserie, a locuit cu ei și lucra, căci de meserie erau constructori de corturi.

\VS{4}În fiecare sâmbătă, el discuta în sinagogă și convingea iudei și greci.


\par }{\PP \VS{5}Când Sila și Timotei s-au coborât din Macedonia, Pavel a fost împins de Duhul Sfânt să mărturisească iudeilor că Isus este Hristosul.

\VS{6}Când aceștia i s-au împotrivit și au blasfemiat, el și-a scuturat hainele și le-a zis: “Sângele vostru să fie pe capetele voastre! Eu sunt curat! De acum înainte, voi merge la neamuri!”.


\par }{\PP \VS{7}A plecat de acolo și a intrat în casa unui om numit Iustus, care se închina lui Dumnezeu și a cărui casă era lângă sinagogă.

\VS{8}Crispus, conducătorul sinagogii, a crezut în Domnul cu toată casa lui. Mulți dintre corinteni, când au auzit, au crezut și au fost botezați.

\VS{9}Domnul i-a spus lui Pavel noaptea, printr-o viziune:
{\WJ{“Nu te teme, ci vorbește și nu tăcea;
}}

\VS{10}{\WJ{căci Eu sunt cu tine și nimeni nu te va ataca ca să-ți facă rău, căci am mulți oameni în această cetate.”
}}


\par }{\PP \VS{11}Și a locuit acolo un an și șase luni, învățând cuvântul lui Dumnezeu printre ei.

\VS{12}Dar, când Gallio era proconsul al Ahaiei, iudeii s-au ridicat de comun acord împotriva lui Pavel și l-au adus înaintea scaunului de judecată,

\VS{13}spunând: “Omul acesta convinge pe oameni să se închine lui Dumnezeu contrar Legii”.


\par }{\PP \VS{14}Dar când Pavel era pe punctul de a deschide gura, Gallio a zis iudeilor: “Dacă ar fi vorba de o greșeală sau de o crimă rea, iudeilor, ar fi cuviincios să vă îngădui;

\VS{15}dar dacă este vorba de cuvinte și de nume și de legea voastră, căutați voi înșivă. Căci eu nu vreau să fiu judecător în aceste chestiuni”.

\VS{16}Și i-a alungat de pe scaunul de judecată.


\par }{\PP \VS{17}Atunci toți grecii au prins pe Sostene, căpetenia sinagogii, și l-au bătut înaintea scaunului de judecată. Lui Gallio nu i-a păsat de niciunul dintre aceste lucruri.


\par }{\PP \VS{18}Pavel, după ce a mai stat încă multe zile, a luat rămas bun de la frați
\FTNT{*}{{\FR 18:18 }{\FT{Cuvântul “frați” aici și acolo unde contextul permite poate fi tradus corect și prin “frați și surori” sau “frați”.}}}și a plecat de acolo în Siria, împreună cu Priscila și Aquila. La Cencreea și-a ras capul, căci avea un jurământ.

\VS{19}A ajuns la Efes și i-a lăsat acolo; dar el însuși a intrat în sinagogă și a discutat cu iudeii.

\VS{20}Când aceștia l-au rugat să rămână mai mult timp cu ei, el a refuzat;

\VS{21}dar, luându-și rămas bun de la ei, a spus: “Trebuie să țin cu orice preț sărbătoarea care urmează la Ierusalim, dar mă voi întoarce la voi dacă Dumnezeu va vrea”. Apoi a pornit din Efes.


\par }{\PP \VS{22}După ce a debarcat la Cezareea, s-a suit, a salutat adunarea și s-a coborât la Antiohia.

\VS{23}După ce a petrecut ceva timp acolo, a plecat și a străbătut în ordine regiunea Galatiei și Frigiei, stabilind pe toți discipolii.

\VS{24}Un iudeu pe nume Apolo, de neam alexandrin, om elocvent, a venit la Efes. El era puternic în Scripturi.

\VS{25}Acest om fusese instruit în calea Domnului; și, fiind plin de ardoare în duh, vorbea și învăța cu exactitate lucrurile despre Isus, deși nu cunoștea decât botezul lui Ioan.

\VS{26}A început să vorbească cu îndrăzneală în sinagogă. Dar Priscila și Aquila, când l-au auzit, l-au luat deoparte și i-au explicat mai exact calea lui Dumnezeu.


\par }{\PP \VS{27}După ce s-a hotărât să treacă în Ahaia, frații l-au încurajat și au scris ucenicilor să îl primească. După ce a venit, i-a ajutat foarte mult pe cei care crezuseră prin har;

\VS{28}pentru că i-a combătut cu putere pe iudei, arătând public, prin Scripturi, că Isus este Hristosul.



\par }\Chap{19}{\PP \VerseOne{1}Pe când Apolo era la Corint, Pavel, trecând prin ținutul de sus, a venit la Efes și a găsit niște ucenici.

\VS{2}El i-a întrebat: “Ați primit Duhul Sfânt atunci când ați crezut?”

\par }{\PP Ei i-au spus: “Nu, noi nici măcar nu am auzit că există un Duh Sfânt”.


\par }{\PP \VS{3}El a zis: “În ce ați fost botezați?”

\par }{\PP Ei au spus: “În botezul lui Ioan”.


\par }{\PP \VS{4}Pavel a spus: “Ioan a botezat cu botezul pocăinței, spunând poporului că trebuie să creadă în cel ce va veni după el, adică în Hristos Isus.”
\FTNT{*}{{\FR 19:4 }{\FT{NU îl omite pe Hristos.}}}


\par }{\PP \VS{5}Când au auzit acestea, s-au botezat în Numele Domnului Isus.

\VS{6}După ce Pavel și-a pus mâinile peste ei, Duhul Sfânt s-a pogorât peste ei și au vorbit în alte limbi și au proorocit.

\VS{7}Erau în total vreo doisprezece bărbați.


\par }{\PP \VS{8}Și, intrând în sinagogă, a vorbit cu îndrăzneală timp de trei luni de zile, și a vorbit și a convins despre lucrurile legate de Împărăția lui Dumnezeu.


\par }{\PP \VS{9}Dar, când unii s-au împietrit și au rămas neascultători, vorbind de rău de Calea Domnului înaintea mulțimii, s-a depărtat de ei și a despărțit pe ucenici, învățând în fiecare zi în școala lui Tiranus.

\VS{10}Acest lucru a continuat timp de doi ani, astfel încât toți cei care locuiau în Asia au auzit cuvântul Domnului Isus, atât iudeii, cât și grecii.


\par }{\PP \VS{11}Dumnezeu făcea minuni deosebite prin mâinile lui Pavel,

\VS{12}astfel că până și batiste sau șorțuri erau duse de pe trupul lui la cei bolnavi, și bolile se depărtau de ei, iar duhurile rele ieșeau.

\VS{13}Dar unii dintre iudeii ambulanți, exorciști, se apucau să invoce asupra celor care aveau duhuri rele numele Domnului Isus, spunând: “Vă conjurăm pe Isus pe care îl predică Pavel”.

\VS{14}Erau șapte fii ai unui anume Sceva, un preot-șef iudeu, care făceau acest lucru.


\par }{\PP \VS{15}Duhul cel rău a răspuns: “Pe Isus îl cunosc, pe Pavel îl cunosc, dar tu cine ești?”

\VS{16}Bărbatul în care era duhul rău a sărit asupra lor, i-a copleșit și i-a biruit, așa că au fugit din casa aceea goi și răniți.

\VS{17}Acest lucru a ajuns la cunoștința tuturor, atât iudei cât și greci, care locuiau în Efes. Frica a căzut peste toți și numele Domnului Isus a fost mărit.

\VS{18}Și mulți dintre cei care crezuseră au venit, mărturisind și mărturisind faptele lor.

\VS{19}Mulți dintre cei care practicau artele magice și-au adunat cărțile și le-au ars în văzul tuturor. Au numărat prețul lor și au găsit că era de cincizeci de mii de arginți.
\FTNT{†}{{\FR 19:19 }{\FT{Cei 50.000 de arginți de aici se refereau probabil la 50.000 de drahme. Dacă este așa, valoarea cărților arse era echivalentă cu aproximativ 160 de ani-om de salariu pentru muncitorii agricoli}}}

\VS{20}Astfel, cuvântul Domnului creștea și devenea puternic.


\par }{\PP \VS{21}După ce au trecut aceste lucruri, Pavel a hotărât în Duhul Sfânt, după ce a trecut prin Macedonia și Ahaia, să se ducă la Ierusalim, zicând: “După ce am fost acolo, trebuie să văd și Roma.”


\par }{\PP \VS{22}După ce a trimis în Macedonia pe doi dintre cei ce-i slujeau, Timotei și Erast, a rămas el însuși în Asia pentru o vreme.

\VS{23}În vremea aceea, s-a iscat o tulburare nu mică în legătură cu Calea.

\VS{24}Căci un om pe nume Demetrius, un argintar care făcea sanctuare de argint pentru Artemis, a adus o afacere nu mică meșterilor

\VS{25}pe care i-a adunat împreună cu lucrătorii cu aceeași ocupație și le-a spus: “Domnilor, știți că din această afacere ne avem averea.

\VS{26}Vedeți și auziți că nu numai în Efes, ci aproape în toată Asia, acest Pavel a convins și a îndepărtat pe mulți oameni, spunând că nu sunt dumnezei făcuți de mâini.

\VS{27}Nu numai că există pericolul ca acest comerț al nostru să ajungă în dizgrație, ci și ca templul marii zeițe Artemis să fie socotit ca nimic și să fie distrusă măreția ei, căreia i se închină toată Asia și toată lumea.”


\par }{\PP \VS{28}Când au auzit acestea, s-au umplut de mânie și au strigat: “Mare este Artemis a efesenilor!”.

\VS{29}Toată cetatea s-a umplut de confuzie și s-au năpustit deodată în teatru, după ce i-au prins pe Gaius și Aristarh, bărbați din Macedonia, tovarășii de drum ai lui Pavel.

\VS{30}Când Pavel a vrut să intre în mijlocul poporului, ucenicii nu i-au dat voie.

\VS{31}De asemenea, unii dintre asiarhi, care îi erau prieteni, au trimis la el și l-au rugat să nu se aventureze în teatru.

\VS{32}Unii, așadar, strigau una și alții alta, căci adunarea era în confuzie. Cei mai mulți dintre ei nu știau de ce se adunaseră.

\VS{33}Au adus pe Alexandru din mulțime, Iudeii punându-l în față. Alexandru făcea semn cu mâna și ar fi vrut să ia apărarea poporului.

\VS{34}Dar, când și-au dat seama că era iudeu, toți, cu un singur glas, timp de vreo două ore, au strigat: “Mare este Artemis din Efes!”.


\par }{\PP \VS{35}După ce a liniștit mulțimea, secretarul orașului a zis: “Bărbați din Efes, cine este acela care nu știe că cetatea efesenilor este templu al marii zeițe Artemis și al icoanei căzute de la Zeus?

\VS{36}Văzând deci că aceste lucruri nu pot fi negate, ar trebui să vă liniștiți și să nu faceți nimic nechibzuit.

\VS{37}Căci voi i-ați adus aici pe acești oameni, care nu sunt nici hoți de temple, nici blasfemiatori ai zeiței voastre.

\VS{38}Așadar, dacă Demetrius și meșterii care sunt cu el au ceva împotriva cuiva, tribunalele sunt deschise și sunt proconsuli. Să depună plângere unii împotriva altora.

\VS{39}Dar dacă căutați ceva despre alte chestiuni, se va rezolva în adunarea obișnuită.

\VS{40}Căci, într-adevăr, riscăm să fim acuzați cu privire la revolta de astăzi, fără să existe niciun motiv. Cu privire la aceasta, nu am fi în stare să dăm socoteală de această agitație.”

\VS{41}După ce a vorbit astfel, a concediat adunarea.



\par }\Chap{20}{\PP \VerseOne{1}După ce a încetat zarva, Pavel a chemat pe ucenici, a luat rămas bun de la ei și a plecat să se ducă în Macedonia.

\VS{2}După ce a străbătut acele părți și i-a încurajat cu multe cuvinte, a ajuns în Grecia.

\VS{3}După ce a petrecut acolo trei luni și după ce iudeii au pus la cale un complot împotriva lui, pe când era pe punctul de a pleca în Siria, a hotărât să se întoarcă prin Macedonia.

\VS{4}Aceștia l-au însoțit până în Asia: Sopater din Beroea, Aristarchus și Secundus din Tesalonic, Gaius din Derbe, Timotei, precum și Tychicus și Trofimus din Asia.

\VS{5}Dar aceștia plecaseră înainte și ne așteptau la Troa.

\VS{6}Am plecat din Filipi după zilele Azimilor și am ajuns la ei la Troa în cinci zile, unde am stat șapte zile.


\par }{\PP \VS{7}În prima zi a săptămânii, când ucenicii s-au adunat să frângă pâinea, Pavel a vorbit cu ei, cu gândul de a pleca a doua zi, și a continuat să vorbească până la miezul nopții.

\VS{8}Erau multe lumini în camera de sus unde
\FTNT{*}{{\FR 20:8 }{\FT{TR citește “ei” în loc de “noi”}}}eram adunați.

\VS{9}Un tânăr pe nume Eutihie ședea la fereastră, împovărat de un somn adânc. Pe când Pavel vorbea încă și mai mult, fiind îngreunat de somn, a căzut de la etajul al treilea și a fost luat mort.

\VS{10}Pavel s-a coborât, a căzut peste el și, îmbrățișându-l, i-a zis: “Nu te tulbura, căci viața lui este în el.”


\par }{\PP \VS{11}După ce S-a suit, a frânt pâinea și a mâncat, și a stat de vorbă cu ei multă vreme, până la ziuă, a plecat.

\VS{12}Ei au adus băiatul viu și au fost foarte mângâiați.


\par }{\PP \VS{13}Dar noi, mergând înainte la corabie, am pornit spre Assos, ca să îmbarcăm acolo pe Pavel, căci el însuși voia să meargă pe uscat.

\VS{14}Când ne-a întâlnit la Assos, l-am luat la bord și am ajuns la Mitilene.

\VS{15}Navigând de acolo, am ajuns a doua zi în fața Chiosului. A doua zi am atins Samos și am rămas la Trogyllium, iar a doua zi am ajuns la Milet.

\VS{16}Căci Pavel hotărâse să treacă pe lângă Efes, ca să nu fie nevoit să petreacă timp în Asia, pentru că se grăbea, dacă îi era posibil, să fie la Ierusalim în ziua Cincizecimii.


\par }{\PP \VS{17}Din Milet a trimis la Efes și a chemat la el pe bătrânii adunării.

\VS{18}După ce au venit la el, le-a zis: “Voi înșivă știți, din prima zi în care am pus piciorul în Asia, cum am fost tot timpul cu voi,

\VS{19}slujindu-l pe Domnul cu toată smerenia, cu multe lacrimi și cu încercările care mi se întâmplau din cauza uneltirilor iudeilor;

\VS{20}cum nu m-am sfiit să vă spun tot ce era de folos, învățându-vă în public și din casă în casă,

\VS{21}mărturisind atât iudeilor, cât și grecilor, pocăința față de Dumnezeu și credința în Domnul nostru Isus.
\FTNT{†}{{\FR 20:21 }{\FT{TR adaugă “Hristos”}}}

\VS{22}Acum, iată, mă duc legat de Duhul Sfânt la Ierusalim, fără să știu ce mi se va întâmpla acolo;

\VS{23}doar că Duhul Sfânt mărturisește în fiecare cetate, spunând că mă așteaptă legături și necazuri.

\VS{24}Dar aceste lucruri nu contează, și nici nu țin la viața mea, ca să îmi termin cu bucurie cursa și slujba pe care am primit-o de la Domnul Isus, ca să mărturisesc pe deplin Vestea cea bună a harului lui Dumnezeu.


\par }{\PP \VS{25}“Acum, iată, știu că voi toți, printre care am umblat propovăduind Împărăția lui Dumnezeu, nu-mi veți mai vedea fața.

\VS{26}De aceea vă mărturisesc astăzi că sunt curat de sângele tuturor oamenilor,

\VS{27}căci nu m-am sfiit să vă vestesc tot sfatul lui Dumnezeu.

\VS{28}Luați seama, așadar, la voi înșivă și la toată turma, în care Duhul Sfânt v-a pus supraveghetori, ca să păstoriți adunarea Domnului și\FTNT{‡}{{\FR 20:28 }{\FT{TR, NU omit “Domnul și”}}} Dumnezeului, pe care a cumpărat-o cu propriul sânge.

\VS{29}Căci știu că, după plecarea mea, vor intra în mijlocul vostru lupi vicioși, care nu vor cruța turma.

\VS{30}Se vor ridica oameni din mijlocul vostru, care vor vorbi lucruri perverse, ca să atragă pe ucenici după ei.

\VS{31}De aceea, vegheați, aducându-vă aminte că, timp de trei ani, nu am încetat să avertizez pe toată lumea zi și noapte cu lacrimi.

\VS{32}Acum, fraților, vă încredințez lui Dumnezeu și cuvântului harului său, care este în măsură să vă zidească și să vă dea moștenirea între toți cei sfințiți.

\VS{33}Nu am râvnit la argintul, aurul sau îmbrăcămintea nimănui.

\VS{34}Voi înșivă știți că aceste mâini au slujit nevoilor mele și ale celor care erau cu mine.

\VS{35}În toate v-am dat ca exemplu că, lucrând astfel, trebuie să-i ajutați pe cei slabi și să vă aduceți aminte de cuvintele Domnului Isus, care a spus el însuși: “{\WJ{Este mai ferice să dai decât să primești”.”
}}


\par }{\PP \VS{36}După ce a spus aceste lucruri, a îngenuncheat și s-a rugat împreună cu toți.

\VS{37}Toți plângeau în hohote de plâns, se aruncau la gâtul lui Pavel și-l sărutau,

\VS{38}maimult decât orice altceva, întristați de cuvântul pe care-l spusese, că nu-i vor mai vedea fața. Apoi l-au însoțit până la corabie.



\par }\Chap{21}{\PP \VerseOne{1}După ce ne-am despărțit de ei și am pornit, am ajuns în linie dreaptă la Cos, iar a doua zi la Rodos și de acolo la Patara.

\VS{2}După ce am găsit o corabie care traversa spre Fenicia, ne-am urcat la bord și am pornit.

\VS{3}După ce am ajuns în dreptul Ciprului, lăsându-l la stânga, am navigat spre Siria și am debarcat la Tir, căci nava se afla acolo pentru a-și descărca încărcătura.

\VS{4}După ce am găsit discipoli, am rămas acolo șapte zile. Aceștia i-au spus lui Pavel, prin Duhul Sfânt, că nu trebuie să se urce la Ierusalim.

\VS{5}După ce au trecut acele zile, am plecat și ne-am continuat călătoria. Toți, cu soțiile și copiii, ne-au însoțit pe drum până când am ieșit din oraș. Îngenunchind pe plajă, ne-am rugat.

\VS{6}După ce ne-am luat rămas bun, ne-am urcat pe corabie, iar ei s-au întors din nou acasă.


\par }{\PP \VS{7}După ce am plecat din Tir, am ajuns la Ptolemaida. I-am întâmpinat pe frați și am rămas cu ei o zi.

\VS{8}A doua zi, noi, care eram tovarășii lui Pavel, am plecat și am ajuns la Cezareea.

\par }{\PP Am intrat în casa lui Filip, evanghelistul, care era unul dintre cei șapte, și am rămas cu el.

\VS{9}Omul acesta avea patru fiice fecioare care profețeau.

\VS{10}Pe când stăteam acolo câteva zile, a coborât din Iudeea un profet numit Agabus.

\VS{11}Venind la noi și luând centura lui Pavel, și-a legat singur picioarele și mâinile și a zis: “Duhul Sfânt spune: “Astfel, iudeii din Ierusalim vor lega pe omul care are această centură și îl vor da în mâinile neamurilor.””


\par }{\PP \VS{12}Când am auzit aceste lucruri, noi și oamenii din locul acela l-am rugat să nu se suie la Ierusalim.

\VS{13}Atunci Pavel a răspuns: “Ce faceți, plângeți și îmi frângeți inima? Căci sunt gata nu numai să fiu legat, ci și să mor la Ierusalim pentru Numele Domnului Isus.”


\par }{\PP \VS{14}Și cum nu se lăsa convins, am încetat, zicând: “Facă-se voia Domnului!”


\par }{\PP \VS{15}După aceste zile, ne-am luat bagajele și ne-am suit la Ierusalim.

\VS{16}Unii dintre ucenicii din Cezareea au mers și ei cu noi, aducând și pe un anume Mnason din Cipru, un ucenic dintâi, la care urma să rămânem.


\par }{\PP \VS{17}Când am ajuns la Ierusalim, frații ne-au primit cu bucurie.

\VS{18}A doua zi, Pavel a intrat cu noi la Iacov; și toți bătrânii erau de față.

\VS{19}După ce i-a salutat, a relatat unul câte unul lucrurile pe care le făcuse Dumnezeu printre neamuri prin slujba lui.

\VS{20}Aceștia, când au auzit, au slăvit pe Dumnezeu. I-au zis: “Vezi, frate, câte mii sunt printre iudei dintre cei care au crezut și toți sunt zeloși pentru Lege.

\VS{21}Ei au fost informați despre tine, că îi înveți pe toți iudeii care sunt printre neamuri să se lepede de Moise, spunându-le să nu-și taie împrejur copiii și să nu umble după obiceiuri.

\VS{22}Și atunci, ce este? Adunarea trebuie să se întrunească cu siguranță, căci vor auzi că ai venit.

\VS{23}Faceți deci ceea ce vă spunem. Avem patru bărbați care au făcut un jurământ.

\VS{24}Ia-i și purifică-te împreună cu ei și plătește-le cheltuielile pentru ei, ca să-și radă capul. Atunci vor ști cu toții că nu este niciun adevăr în lucrurile despre care au fost informați despre tine, ci că și tu însuți umbli respectând Legea.

\VS{25}Dar în ceea ce privește neamurile care cred, am scris hotărârea noastră că nu trebuie să respecte nimic de acest fel, decât să se ferească de mâncarea oferită idolilor, de sânge, de lucrurile sugrumate și de imoralitatea sexuală.”


\par }{\PP \VS{26}Pavel a luat pe bărbați, s-a curățit a doua zi și a mers cu ei în Templu, și a vestit împlinirea zilelor de curățire, până când s-a adus jertfă pentru fiecare dintre ei.

\VS{27}Când cele șapte zile erau aproape împlinite, iudeii din Asia, văzându-l în templu, au stârnit toată mulțimea și au pus mâinile pe el,

\VS{28}strigând: “Bărbați ai lui Israel, ajutor! Acesta este omul care îi învață pe toți oamenii de pretutindeni împotriva poporului, a legii și a acestui loc. Mai mult, el a adus și greci în templu și a pângărit acest loc sfânt!”

\VS{29}Căci îl văzuseră pe Trofim, efescianul, împreună cu el în cetate, și au presupus că Pavel l-a adus în templu.


\par }{\PP \VS{30}Toată cetatea s-a mișcat și poporul a alergat. L-au prins pe Pavel și l-au târât afară din templu. Imediat ușile au fost închise.

\VS{31}În timp ce încercau să-l ucidă, a sosit la comandantul regimentului vestea că tot Ierusalimul era în zarvă.

\VS{32}Imediat a luat soldați și centurioni și a alergat la ei. Aceștia, când i-au văzut pe căpitanul principal și pe soldați, au încetat să-l mai bată pe Pavel.

\VS{33}Atunci comandantul s-a apropiat, l-a arestat, a poruncit să fie legat cu două lanțuri și a întrebat cine este și ce a făcut.

\VS{34}Unii strigau un lucru și alții altul, în mulțime. Când nu a putut afla adevărul din cauza gălăgiei, a poruncit să fie dus în cazarmă.


\par }{\PP \VS{35}Când a ajuns pe scări, a fost dus de ostași, din pricina mulțimii care se înghesuia;

\VS{36}căci mulțimea poporului îl urmărea, strigând: “Luați-l!”.

\VS{37}Pe când Pavel era pe punctul de a fi dus în cazarmă, l-a întrebat pe comandant: “Pot să-ți vorbesc?”

\par }{\PP El a spus: “Știi greacă?

\VS{38}Nu ești tu, deci, egipteanul care, mai înainte de aceste zile, a stârnit răzvrătirea și a condus în pustiu pe cei patru mii de oameni ai asasinilor?”


\par }{\PP \VS{39}Dar Pavel a zis: “Eu sunt un iudeu din Tars, din Cilicia, un cetățean al unei cetăți deloc neînsemnate. Vă rog, dați-mi voie să vorbesc poporului.”


\par }{\PP \VS{40}După ce i-a dat voie, Pavel, stând pe scări, a făcut semn cu mâna către popor. Când s-a făcut o mare tăcere, el le-a vorbit în limba ebraică, zicând



\par }\Chap{22}{\PP \VerseOne{1}“Frați și părinți, ascultați apărarea pe care v-o fac acum.”


\par }{\PP \VS{2}Când au auzit că le vorbea în limba ebraică, s-au liniștit și mai mult.

\par }{\PP El a zis:

\VS{3}“Eu sunt iudeu, născut în Tarsul Ciliciei, dar crescut în această cetate, la picioarele lui Gamaliel, învățat după tradiția strictă a legii părinților noștri, și zelos pentru Dumnezeu, așa cum sunteți și voi toți astăzi.

\VS{4}Am prigonit această Cale până la moarte, legând și dând în închisori bărbați și femei,

\VS{5}așa cum mărturisesc și marele preot și tot consiliul bătrânilor, de la care am primit și scrisori către frați, și am călătorit la Damasc pentru a-i aduce în legături la Ierusalim și pe cei care erau acolo, ca să fie pedepsiți.


\par }{\PP \VS{6}“Pe când mergeam și mă apropiam de Damasc, pe la amiază, deodată, din cer, o lumină mare a strălucit în jurul meu.

\VS{7}Am căzut la pământ și am auzit un glas care îmi spunea:
{\WJ{“Saul, Saul, de ce mă prigonești?”
}}

\VS{8}I-am răspuns: “Cine ești Tu, Doamne?” El mi-a zis: “{\WJ{Eu sunt Isus din Nazaret, pe care Îl prigonești tu.
}}


\par }{\PP \VS{9}Cei ce erau cu mine au văzut lumina și s-au temut, dar n-au înțeles vocea Celui ce vorbea cu mine.

\VS{10}Am întrebat: “Ce să fac, Doamne?”. Domnul mi-a zis:
{\WJ{'Ridică-te și du-te în Damasc. Acolo ți se vor spune toate lucrurile care ți-au fost rânduite să le faci”.
}}

\VS{11}Când nu mai vedeam nimic din cauza slavei acelei lumini, fiind condus de mâna celor care erau cu mine, am intrat în Damasc.


\par }{\PP \VS{12}“Un oarecare Anania, un om evlavios după Lege, vestit de toți iudeii care locuiau în Damasc,

\VS{13}a venit la mine și, stând lângă mine, mi-a zis: “Frate Saul, primește-ți vederea! Chiar în acel ceas m-am uitat la el.

\VS{14}El a zis: “Dumnezeul părinților noștri te-a rânduit să cunoști voia Lui, să-L vezi pe Cel Drept și să auzi un glas din gura Lui.

\VS{15}Căci tu vei fi martor pentru el, în fața tuturor oamenilor, pentru ceea ce ai văzut și ai auzit.

\VS{16}Și acum, de ce așteptați? Ridică-te, botează-te și spală-ți păcatele, invocând numele Domnului”.


\par }{\PP \VS{17}După ce m-am întors la Ierusalim și în timp ce mă rugam în Templu, am căzut în transă

\VS{18}și l-am văzut pe El zicându-mi: “{\WJ{Grăbește-te și ieși repede din Ierusalim, pentru că nu vor primi de la tine mărturie despre Mine”.
}}

\VS{19}Eu am zis: “Doamne, ei înșiși știu că am întemnițat și am bătut în toate sinagogile pe cei care credeau în Tine.

\VS{20}Când s-a vărsat sângele lui Ștefan, martorul Tău, am stat și eu de față, consimțind la moartea lui și păzind mantiile celor care l-au ucis'.


\par }{\PP \VS{21}El mi-a zis: “{\WJ{Pleacă, căci te voi trimite departe de aici, la neamuri.”
}}


\par }{\PP \VS{22}L-au ascultat până ce a spus aceasta; apoi au ridicat glasul și au zis: “Scapă pământul de omul acesta, căci nu este vrednic să trăiască!”


\par }{\PP \VS{23}Pe când strigau, își aruncau mantiile și aruncau praf în aer,

\VS{24}comandantul a poruncit să fie dus în cazarmă și a ordonat să fie cercetat prin biciuire, ca să știe pentru ce crimă strigau așa împotriva lui.

\VS{25}După ce l-au legat cu chingi, Pavel l-a întrebat pe centurionul care se afla de față: “Îți este permis să biciuiești un om care este roman și nu a fost găsit vinovat?”


\par }{\PP \VS{26}Când a auzit centurionul, s-a dus la comandant și i-a zis: “Ai grijă ce ai de gând să faci, căci omul acesta este roman!”


\par }{\PP \VS{27}Comandantul a venit și l-a întrebat: “Spune-mi, ești roman?”

\par }{\PP El a spus: “Da”.


\par }{\PP \VS{28}Și comandantul a răspuns: “Mi-am cumpărat cetățenia cu un preț mare.”

\par }{\PP Pavel a spus: “Dar eu m-am născut roman”.


\par }{\PP \VS{29}Și îndată s-au depărtat de el cei ce voiau să-l cerceteze, iar comandantul s-a temut și el, când a văzut că este roman, pentru că îl legase.

\VS{30}Dar a doua zi, dorind să afle adevărul despre motivul pentru care era acuzat de iudei, l-a eliberat din legături și a poruncit preoților de seamă și întregului consiliu să se adune, a coborât pe Pavel și l-a pus în fața lor.



\par }\Chap{23}{\PP \VerseOne{1}Pavel, uitându-se cu stăruință la consiliu, a zis: “Fraților, până astăzi am trăit înaintea lui Dumnezeu cu toată conștiința împăcată”.


\par }{\PP \VS{2}Marele preot Anania a poruncit celor ce stăteau lângă el să îl lovească peste gură.


\par }{\PP \VS{3}Atunci Pavel i-a zis: “Dumnezeu te va lovi, zid albit! Stai tu să mă judeci după lege și poruncești să fiu lovit contrar legii?”


\par }{\PP \VS{4}Cei ce stăteau de față ziceau: “Voi defăimați pe marele preot al lui Dumnezeu?”


\par }{\PP \VS{5}Pavel a zis: “Nu știam, fraților, că este mare preot. Căci este scris: “Să nu vorbești de rău despre un conducător al poporului tău”.”
\FTNT{*}{{\FR 23:5 }{\FT{Exodul 22:28}}}


\par }{\PP \VS{6}Dar Pavel, văzând că o parte erau saduchei, iar cealaltă farisei, a strigat în consiliu: “Bărbați și frați, eu sunt fariseu, fiu de farisei. Cu privire la speranța și la învierea morților sunt judecat!”


\par }{\PP \VS{7}După ce a spus acestea, s-a iscat o ceartă între farisei și saduchei, și mulțimea s-a dezbinat.

\VS{8}Căci saducheii spuneau că nu există nici înviere, nici înger, nici duh; dar fariseii mărturisesc toate acestea.

\VS{9}S-a stârnit o mare zarvă și unii dintre cărturarii din partea fariseilor s-au ridicat în picioare și se certau, zicând: “Noi nu găsim nimic rău în omul acesta. Dar dacă un duh sau un înger i-a vorbit, să nu ne luptăm împotriva lui Dumnezeu!”


\par }{\PP \VS{10}Când s-a iscat o mare ceartă, comandantul, temându-se ca Pavel să nu fie sfâșiat de ei, a poruncit ostașilor să se coboare, să-l ia cu forța dintre ei și să-l ducă în cazarmă.


\par }{\PP \VS{11}În noaptea următoare, Domnul a stat lângă el și i-a zis:
{\WJ{“Înveselește-te, Pavel, căci, după cum ai mărturisit despre Mine la Ierusalim, tot așa trebuie să mărturisești și la Roma.”
}}


\par }{\PP \VS{12}Când s-a făcut ziuă, unii dintre iudei s-au unit și s-au legat cu un blestem, zicând că nu vor mânca și nu vor bea nimic până ce nu vor omorî pe Pavel.

\VS{13}Erau mai mult de patruzeci de persoane care făcuseră această conspirație.

\VS{14}Ei au venit la preoții cei mai de seamă și la bătrâni și au zis: “Ne-am legat sub un mare blestem să nu gustăm nimic până nu-l vom ucide pe Pavel.

\VS{15}Acum, deci, voi, împreună cu consiliul, informați comandantul că mâine trebuie să vi-l aducă la voi, ca și cum ați avea de gând să judecați mai exact cazul lui. Noi suntem gata să-l omorâm înainte de a se apropia.”


\par }{\PP \VS{16}Dar fiul surorii lui Pavel a auzit că stau la pândă, a venit și a intrat în cazarmă și a anunțat pe Pavel.

\VS{17}Pavel a chemat pe unul dintre centurioni și a zis: “Duceți-l pe acest tânăr la comandant, căci are ceva să-i spună.”


\par }{\PP \VS{18}L-a luat, l-a dus la comandant și i-a zis: “Pavel, prizonierul, m-a chemat și m-a rugat să-ți aduc pe acest tânăr. El are ceva să-ți spună”.


\par }{\PP \VS{19}Comandantul l-a luat de mână și, ducându-se la o parte, l-a întrebat în particular: “Ce ai să-mi spui?”


\par }{\PP \VS{20}Și a zis: “Iudeii au căzut de acord să te roage să aduci mâine pe Pavel la sfat, ca și cum ar fi vrut să se intereseze mai bine de el.

\VS{21}De aceea, nu le ceda, pentru că îl pândesc mai mult de patruzeci de oameni care s-au legat cu un blestem să nu mănânce și să nu bea nimic până nu-l vor ucide. Acum sunt gata, așteptând promisiunea din partea ta.”


\par }{\PP \VS{22}Și comandantul a lăsat pe tânăr să plece, și i-a zis: “Să nu spui nimănui că mi-ai descoperit aceste lucruri.”


\par }{\PP \VS{23}A chemat la el doi dintre centurioni și le-a zis: “Pregătiți două sute de ostași, ca să meargă până la Cezareea, cu șaptezeci de călăreți și două sute de oameni înarmați cu sulițe, la ceasul al treilea din noapte”.
\FTNT{†}{{\FR 23:23 }{\FT{în jurul orei 21:00.}}}

\VS{24}Le-a cerut să pună la dispoziție niște călăreți, ca să-l urce pe Pavel pe unul dintre ei și să-l ducă în siguranță la Felix, guvernatorul.

\VS{25}Și a scris o scrisoare astfel: “Și a scris o scrisoare ca aceasta:


\par }{\PP \VS{26}“Claudius Lysias către cel mai bun guvernator Felix: “Salutări.


\par }{\PP \VS{27}“Omul acesta a fost prins de iudei și era pe cale să fie omorât de ei, când am venit cu ostașii și l-am salvat, aflând că este roman.

\VS{28}Dorind să aflu cauza pentru care îl acuzau, l-am adus jos la consiliul lor.

\VS{29}Am găsit că era acuzat în legătură cu chestiuni legate de legea lor, dar nu era acuzat de nimic vrednic de moarte sau de închisoare.

\VS{30}Când mi s-a spus că iudeii îl pândesc, l-am trimis imediat la voi, însărcinându-i și pe acuzatorii lui să aducă înaintea voastră acuzațiile lor împotriva lui. Rămas bun!”


\par }{\PP \VS{31}Soldații au luat pe Pavel și, după porunca lor, l-au dus noaptea la Antipatris.

\VS{32}Dar a doua zi au lăsat călăreții să meargă cu el și s-au întors la cazarmă.

\VS{33}Când au ajuns la Cezareea și au predat scrisoarea guvernatorului, i-au prezentat și lui Pavel.

\VS{34}După ce a citit-o, guvernatorul a întrebat din ce provincie este. Când a înțeles că era din Cilicia, a spus:

\VS{35}“Te voi asculta pe deplin când vor sosi și acuzatorii tăi”. El a poruncit să fie ținut în palatul lui Irod.



\par }\Chap{24}{\PP \VerseOne{1}După cinci zile, marele preot Anania s-a coborât cu niște bătrâni și cu un orator, Tertul. Aceștia l-au informat pe guvernator împotriva lui Pavel.

\VS{2}Când a fost chemat, Tertullus a început să-l acuze, spunând: “Văzând că prin tine ne bucurăm de multă pace și că prosperitatea vine în acest popor prin previziunea ta,

\VS{3}o acceptăm în toate felurile și în toate locurile, preasfințite Felix, cu toată recunoștința.

\VS{4}Dar, ca să nu te întârzii, te rog să ai răbdare cu noi și să asculți câteva cuvinte.

\VS{5}Pentru că am descoperit că acest om este o plagă, un instigator de insurecții printre toți evreii din întreaga lume și un lider al sectei nazarienilor.

\VS{6}A încercat chiar să profaneze templul și l-am arestat.
\FTNT{*}{{\FR 24:6 }{\FT{TR adaugă: “Am vrut să-l judecăm după legea noastră,”}}}

\VS{7}\FTNT{†}{{\FR 24:7 }{\FT{TR adaugă “dar comandantul, Lisias, a venit și cu mare violență l-a luat din mâinile noastre,”}}}

\VS{8}\FTNT{‡}{{\FR 24:8 }{\FT{TR adaugă “poruncind acuzatorilor săi să vină la tine”.}}} Examinându-l tu însuți, poți constata toate aceste lucruri de care îl acuzăm.”


\par }{\PP \VS{9}Și iudeii s-au alăturat și ei atacului, afirmând că așa stau lucrurile.


\par }{\PP \VS{10}După ce guvernatorul i-a făcut semn să vorbească, Pavel a răspuns: “Fiindcă știu că tu ești de mulți ani judecător al acestui neam, îmi iau cu plăcere apărarea,

\VS{11}fiindcă poți să constați că nu au trecut mai mult de douăsprezece zile de când m-am suit să mă închin la Ierusalim.

\VS{12}În templu nu m-au găsit certându-mă cu nimeni sau stârnind mulțimea, nici în sinagogi, nici în oraș.

\VS{13}Și nici nu vă pot dovedi lucrurile de care mă acuză acum.

\VS{14}Dar vă mărturisesc acest lucru: că, potrivit Căii, pe care ei o numesc sectă, așa slujesc eu Dumnezeului părinților noștri, crezând toate lucrurile care sunt conform Legii și care sunt scrise în profeți;

\VS{15}având speranța față de Dumnezeu, pe care și aceștia o așteaptă, că va fi o înviere a morților, atât a celor drepți, cât și a celor nedrepți.

\VS{16}În acest sens, practic și eu, având întotdeauna o conștiință lipsită de supărare față de Dumnezeu și față de oameni.

\VS{17}După câțiva ani, am venit să aduc națiunii mele daruri pentru cei nevoiași și ofrande;

\VS{18}în mijlocul cărora niște iudei din Asia m-au găsit purificat în templu, nu cu o mulțime și nici cu agitație.

\VS{19}Ei ar fi trebuit să fie aici, înaintea voastră, și să mă acuze, dacă aveau ceva împotriva mea.

\VS{20}Sau, altfel, să spună ei înșiși ce nedreptate au găsit în mine când am stat în fața consiliului,

\VS{21}dacă nu cumva pentru acest singur lucru am strigat stând în picioare în mijlocul lor: “Cu privire la învierea morților sunt judecat astăzi în fața voastră!””.


\par }{\PP \VS{22}Dar Felix, care cunoștea mai bine Calea, i-a amânat, zicând: “Când va coborî Lisias, comandantul, voi lua o hotărâre în privința voastră.”

\VS{23}El a poruncit centurionului ca Pavel să fie ținut în custodie și să aibă unele privilegii și să nu interzică niciunui dintre prietenii săi să îl servească sau să îl viziteze.


\par }{\PP \VS{24}După câteva zile, Felix a venit cu Drusila, soția sa, care era iudeică, și a trimis să cheme pe Pavel și l-a ascultat cu privire la credința în Hristos Isus.

\VS{25}În timp ce el vorbea despre neprihănire, despre stăpânirea de sine și despre judecata viitoare, Felix s-a îngrozit și i-a răspuns: “Du-te deocamdată și, când îmi va fi de folos, te voi chema.”

\VS{26}Între timp, el mai spera că Pavel îi va da bani, ca să-l elibereze. De aceea, de asemenea, trimitea mai des după el și vorbea cu el.


\par }{\PP \VS{27}După ce s-au împlinit doi ani, Felix a fost înlocuit de Porcius Festus; și Felix, vrând să se facă plăcut iudeilor, a lăsat pe Pavel în temniță.



\par }\Chap{25}{\PP \VerseOne{1}Festus, deci, a venit în provincie și, după trei zile, s-a suit din Cezareea la Ierusalim.

\VS{2}Marele preot și cei mai de seamă dintre iudei l-au informat împotriva lui Pavel și l-au rugat,

\VS{3}cerându-i o favoare împotriva lui, să-l cheme la Ierusalim, complotând să-l ucidă pe drum.

\VS{4}Festus însă a răspuns că Pavel trebuie să fie ținut în custodie la Cezareea și că el însuși urma să plece în curând.

\VS{5}“Lăsați deci”, a spus el, “pe cei care sunt în putere printre voi să coboare cu mine și, dacă este ceva rău în acest om, să îl acuze.”


\par }{\PP \VS{6}După ce a stat printre ei mai mult de zece zile, s-a pogorât la Cezareea și a doua zi a stat pe scaunul de judecată și a poruncit să fie adus Pavel.

\VS{7}După ce a venit, iudeii care se coborâseră de la Ierusalim au stat în jurul lui, aducând împotriva lui multe și grave acuzații pe care nu le puteau dovedi,

\VS{8}în timp ce el spunea în apărarea sa: “Nu am păcătuit deloc nici împotriva legii iudeilor, nici împotriva templului, nici împotriva lui Cezar”.


\par }{\PP \VS{9}Dar Festus, vrând să-și câștige favoarea iudeilor, a luat cuvântul și a zis lui Pavel: “Vrei să te duci la Ierusalim și să fii judecat de mine acolo despre aceste lucruri?”


\par }{\PP \VS{10}Dar Pavel a zis: “Eu stau în fața scaunului de judecată al Cezarului, unde trebuie să fiu judecat. Nu am făcut niciun rău iudeilor, după cum bine știți și voi.

\VS{11}Căci, dacă am greșit și am săvârșit ceva vrednic de moarte, nu refuz să mor; dar dacă nu este adevărat niciunul din lucrurile de care mă acuză ei, nimeni nu mă poate preda lor. Eu apelez la Cezar!”


\par }{\PP \VS{12}Atunci Festus, după ce a vorbit cu consiliul, a răspuns: “Ați apelat la Cezar. La Cezar vă veți duce”.


\par }{\PP \VS{13}După ce au trecut câteva zile, regele Agripa și Berenice au sosit la Cezareea și au salutat pe Festus.

\VS{14}Cum a stat acolo multe zile, Festus a prezentat împăratului cazul lui Pavel, zicând: “Există un om lăsat prizonier de Felix;

\VS{15}despre care, când eram la Ierusalim, m-au informat preoții cei mai de seamă și bătrânii iudeilor, cerând o sentință împotriva lui.

\VS{16}Le-am răspuns că nu este obiceiul romanilor să predea pe cineva la pieire înainte ca acuzatul să se întâlnească față în față cu acuzatorii și să aibă ocazia să se apere în legătură cu acuzațiile care i se aduc.

\VS{17}Prin urmare, când s-au adunat aici, nu am întârziat, ci a doua zi am stat pe scaunul de judecată și am poruncit ca omul să fie adus.

\VS{18}Când acuzatorii s-au ridicat în picioare, nu i-au adus nicio acuzație de lucruri pe care le presupuneam eu,

\VS{19}ci au avut anumite întrebări împotriva lui cu privire la propria lor religie și la un Isus mort, despre care Pavel a afirmat că este viu.

\VS{20}Fiind nedumerit cum să întreb despre aceste lucruri, l-am întrebat dacă este dispus să meargă la Ierusalim și acolo să fie judecat cu privire la aceste lucruri.

\VS{21}Dar, când Pavel a cerut să fie reținut în vederea deciziei împăratului, am poruncit să fie reținut până când îl voi putea trimite la Cezar.”


\par }{\PP \VS{22}Agripa a zis lui Festus: “Aș vrea și eu să aud pe acest om.”

\par }{\PP “Mâine”, a spus el, “îl veți auzi”.


\par }{\PP \VS{23}A doua zi, după ce Agripa și Berenice au venit cu mare pompă și au intrat în sala de audieri, împreună cu comandanții și cu cei mai de seamă din cetate, din porunca lui Festus, a fost adus Pavel.

\VS{24}Festus a zis: “Împărate Agripa, și toți bărbații care sunt aici de față cu noi, vedeți pe acest om despre care toată mulțimea iudeilor mi-a făcut o petiție, atât la Ierusalim, cât și aici, strigând că nu trebuie să mai trăiască.

\VS{25}Dar când am constatat că nu a săvârșit nimic vrednic de moarte și cum el însuși a apelat la împărat, am hotărât să îl trimit pe el,

\VS{26}despre care nu am nimic sigur de scris domnului meu. De aceea l-am adus în fața voastră și mai ales în fața ta, rege Agripa, pentru ca, după examinare, să am ce să scriu.

\VS{27}Căci mi se pare nerezonabil ca, trimițând un prizonier, să nu precizez și acuzațiile care i se aduc.”



\par }\Chap{26}{\PP \VerseOne{1}Agripa a zis lui Pavel: “Poți să vorbești în numele tău.”

\par }{\PP Atunci Pavel a întins mâna și a luat apărarea.

\VS{2}“Mă consider fericit, rege Agripa, că astăzi îmi voi lua apărarea în fața ta cu privire la toate lucrurile de care sunt acuzat de iudei,

\VS{3}mai ales că tu ești expert în toate obiceiurile și în toate chestiunile care există printre iudei. De aceea te rog să mă asculți cu răbdare.


\par }{\PP \VS{4}“Toți iudeii știu cum am trăit din tinerețe, de la început în neamul meu și la Ierusalim,

\VS{5}și m-au cunoscut de la început, dacă vor să mărturisească, că am trăit ca fariseu după cea mai strictă sectă a religiei noastre.

\VS{6}Acum stau aici pentru a fi judecat pentru speranța promisiunii făcute de Dumnezeu părinților noștri,

\VS{7}pe care cele douăsprezece triburi ale noastre, slujind cu sârguință zi și noapte, speră să o obțină. Cu privire la această speranță sunt acuzat de iudei, rege Agripa!

\VS{8}De ce este judecat incredibil la tine dacă Dumnezeu învie morții?


\par }{\PP \VS{9}“Eu însumi am crezut cu adevărat că trebuie să fac multe lucruri potrivnice numelui lui Isus din Nazaret.

\VS{10}Am făcut și eu așa ceva la Ierusalim. Am închis deopotrivă pe mulți dintre sfinți în închisori, primind autoritate de la preoții cei mai de seamă, iar când au fost condamnați la moarte, mi-am dat votul împotriva lor.

\VS{11}Pedepsindu-i adesea în toate sinagogile, am încercat să-i fac să hulească. Fiind extrem de înfuriat împotriva lor, i-am persecutat până în cetăți străine.


\par }{\PP \VS{12}“Pe când călătoream spre Damasc, cu împuternicirea și însărcinarea preoților cei mai de seamă,

\VS{13}la amiază, împărate, am văzut pe drum o lumină din cer, mai strălucitoare decât soarele, care strălucea în jurul meu și al celor ce călătoreau cu mine.

\VS{14}După ce am căzut cu toții la pământ, am auzit o voce care îmi spunea în limba ebraică:
{\WJ{“Saul, Saul, de ce mă prigonești? Îți este greu să lovești cu piciorul în ghiare'.
}}


\par }{\PP \VS{15}Și am zis: “Cine ești Tu, Doamne?

\par }{\PP “El a spus: “{\WJ{Eu sunt Isus, pe care voi Îl persecutați.
}}

\VS{16}{\WJ{Dar ridică-te și stai în picioare, căci pentru aceasta m-am arătat ție: ca să te numesc slujitor și martor atât al lucrurilor pe care le-ai văzut, cât și al celor pe care ți le voi descoperi;
}}

\VS{17}să te
{\WJ{eliberez de la popor și de la neamuri, la care te trimit,
}}

\VS{18}{\WJ{să le deschizi ochii, ca să se întoarcă de la întuneric la lumină și de la puterea lui Satana la Dumnezeu, ca să primească iertarea păcatelor și moștenirea între cei sfințiți prin credința în mine'.
}}


\par }{\PP \VS{19}“De aceea, împărate Agripa, n-am fost neascultător de viziunea cerească,

\VS{20}ci am vestit mai întâi celor din Damasc, la Ierusalim și în toată țara Iudeii, precum și neamurilor, că trebuie să se pocăiască și să se întoarcă la Dumnezeu, făcând fapte vrednice de pocăință.

\VS{21}Din acest motiv, iudeii m-au prins în templu și au încercat să mă ucidă.

\VS{22}De aceea, după ce am obținut ajutorul care vine de la Dumnezeu, stau până în ziua de azi mărturisind atât celor mici, cât și celor mari, fără să spun nimic altceva decât ceea ce profeții și Moise au spus că se va întâmpla,

\VS{23}cum trebuie să sufere Cristos și cum, prin învierea morților, el va fi cel dintâi care va vesti lumina atât acestui popor, cât și neamurilor.”


\par }{\PP \VS{24}Pe când se apăra el astfel, Festus a zis cu glas tare: “Pavel, ești nebun! Marea ta învățătură te face să înnebunești!”


\par }{\PP \VS{25}Dar el a răspuns: “Nu sunt nebun, preafericite Festus, ci spun cu îndrăzneală cuvinte adevărate și rezonabile.

\VS{26}Căci împăratul știe aceste lucruri, căruia îi vorbesc și eu liber. Căci sunt convins că nici unul dintre aceste lucruri nu-i este ascuns, căci acest lucru nu s-a făcut într-un colț.

\VS{27}Împărate Agripa, crezi tu oare în profeți? Știu că tu crezi.”


\par }{\PP \VS{28}Agripa a zis lui Pavel: “Cu puțină convingere vrei să mă faci creștin?”


\par }{\PP \VS{29}Pavel a zis: “Rog pe Dumnezeu ca, fie cu puțin, fie cu mult, nu numai voi, ci și toți cei ce mă ascultă astăzi, să ajungeți la fel ca mine, afară de aceste legături.”


\par }{\PP \VS{30}Împăratul s-a sculat, împreună cu guvernatorul, cu Bernice și cu cei ce stăteau cu ei.

\VS{31}După ce s-au retras, au vorbit unul cu altul, zicând: “Omul acesta nu face nimic vrednic de moarte sau de legături.”

\VS{32}Agrippa i-a spus lui Festus: “Omul acesta ar fi putut fi eliberat dacă nu ar fi apelat la Cezar.”



\par }\Chap{27}{\PP \VerseOne{1}După ce s-a hotărât să ne îmbarcăm spre Italia, au predat pe Pavel și pe alți prizonieri unui centurion numit Iulius, din trupa lui Augustan.

\VS{2}Îmbarcându-ne pe o corabie din Adramyttium, care urma să plece spre locuri de pe coasta Asiei, am pornit pe mare, cu noi fiind Aristarchus, un macedonean din Tesalonic.

\VS{3}A doua zi, am atins Sidonul. Iulius s-a purtat cu Pavel cu bunăvoință și i-a dat voie să se ducă la prietenii săi și să se răcorească.

\VS{4}Plecând de acolo, am navigat sub vântul din Cipru, pentru că vânturile erau contrare.

\VS{5}După ce am traversat marea care se află în largul Ciliciei și al Pamfiliei, am ajuns la Myra, un oraș din Licia.

\VS{6}Acolo, centurionul a găsit o corabie din Alexandria, care naviga spre Italia, și ne-a îmbarcat.

\VS{7}După ce am navigat încet multe zile și am ajuns cu greu în fața Cnidului, vântul nepermițându-ne să mergem mai departe, am navigat sub vântul dinspre Creta, în fața Salmonei.

\VS{8}Navigând cu greu de-a lungul ei, am ajuns la un loc numit Fair Havens, lângă orașul Lasea.


\par }{\PP \VS{9}După ce trecuse mult timp și călătoria era deja periculoasă, pentru că trecuse deja Postul Mare, Pavel i-a sfătuit

\VS{10}și le-a zis: “Domnilor, văd că această călătorie va fi cu pagube și cu multe pierderi, nu numai a încărcăturii și a corăbiei, ci și a vieții noastre.”

\VS{11}Dar centurionul a dat mai multă atenție stăpânului și proprietarului corabiei decât la cele spuse de Pavel.

\VS{12}Pentru că raiul nu era potrivit pentru a ierna, cei mai mulți au sfătuit să plece pe mare de acolo, dacă prin orice mijloace puteau să ajungă la Phoenix și să ierneze acolo, care este un port din Creta, privind spre sud-vest și nord-vest.


\par }{\PP \VS{13}Când vântul de sud a suflat ușor, crezând că și-au atins scopul, au ridicat ancora și au navigat de-a lungul Cretei, aproape de țărm.

\VS{14}Dar, nu după mult timp, un vânt furtunos a bătut dinspre țărm, care se numește Euroclydon.
\FTNT{*}{{\FR 27:14 }{\FT{Sau “o furtună de nord-est”.}}}

\VS{15}Când corabia a fost prinsă și nu a putut face față vântului, am cedat în fața lui și am fost împinși de-a lungul.

\VS{16}Alergând sub vântul unei mici insule numite Clauda, am reușit, cu greu, să asigurăm barca.

\VS{17}După ce au ridicat-o, au folosit cabluri pentru a ajuta la consolidarea navei. Temându-se că vor eșua pe bancurile de nisip din Syrtis, au coborât ancora de mare și astfel au fost împinși de-a lungul.

\VS{18}În timp ce ne chinuiam foarte mult cu furtuna, a doua zi au început să arunce lucruri peste bord.

\VS{19}În a treia zi, au aruncat cu propriile mâini scufundările corăbiei.

\VS{20}Când nici soarele și nici stelele nu ne-au mai luminat timp de multe zile și nici o furtună mică nu ne apăsa, orice speranță că vom fi salvați era acum înlăturată.


\par }{\PP \VS{21}După ce au stat multă vreme fără mâncare, Pavel s-a ridicat în mijlocul lor și le-a zis: “Domnilor, trebuia să mă ascultați și să nu plecați din Creta, ca să nu vă alegeți cu această pagubă și cu această pierdere.

\VS{22}Acum vă îndemn să vă înveseliți, căci nu va fi nicio pierdere de vieți omenești printre voi, ci numai a corabiei.

\VS{23}Căci în această noapte a stat lângă mine un înger, care aparținea Dumnezeului al cărui sunt și căruia îi slujesc,

\VS{24}și care mi-a spus: “Nu te teme, Paul. Trebuie să te prezinți în fața Cezarului. Iată, Dumnezeu ți-a acordat pe toți cei care navighează cu tine.”

\VS{25}Așadar, domnilor, înveseliți-vă! Căci eu cred în Dumnezeu, că va fi așa cum mi s-a spus.

\VS{26}Dar trebuie să eșuăm pe o anumită insulă.”


\par }{\PP \VS{27}Dar, când a venit a paisprezecea noapte, pe când eram împinși de colo-colo în Marea Adriatică, pe la miezul nopții, marinarii au bănuit că se apropiau de un țărm.

\VS{28}Au făcut sondaje și au găsit douăzeci de brațe.\FTNT{†}{{\FR 27:28 }{\FT{20 brazde = 120 picioare = 36,6 metri}}} După puțin timp, au făcut din nou sondaje și au găsit cincisprezece brazde.
\FTNT{‡}{{\FR 27:28 }{\FT{15 brazde = 90 picioare = 27,4 metri}}}

\VS{29}Temându-se că vom eșua pe un teren stâncos, au aruncat patru ancore de la pupa și au dorit să se lumineze de ziuă.

\VS{30}În timp ce marinarii încercau să fugă din corabie și coborâseră barca în mare, prefăcându-se că vor arunca ancorele de la prova,

\VS{31}Pavel le-a spus centurionului și soldaților: “Dacă aceștia nu rămân în corabie, nu puteți fi salvați.”

\VS{32}Atunci soldații au tăiat frânghiile bărcii și au lăsat-o să cadă.


\par }{\PP \VS{33}Pe când se făcea ziuă, Pavel i-a rugat pe toți să ia ceva de mâncare, zicând: “Astăzi este a paisprezecea zi în care așteptați și postiți, fără să luați nimic.

\VS{34}De aceea vă rog să luați ceva de mâncare, căci aceasta este pentru siguranța voastră, căci nu va pieri niciun fir de păr de pe capul niciunuia dintre voi.”

\VS{35}După ce a spus acestea și a luat pâinea, a mulțumit lui Dumnezeu în fața tuturor; apoi a frânt-o și a început să mănânce.

\VS{36}Atunci toți s-au înveselit și au luat și ei mâncare.

\VS{37}În total, eram două sute șaptezeci și șase de suflete pe corabie.

\VS{38}După ce au mâncat destul, au ușurat corabia, aruncând grâul în mare.

\VS{39}Când s-a făcut ziuă, nu au recunoscut pământul, dar au observat un anumit golf cu o plajă și au decis să încerce să conducă corabia pe ea.

\VS{40}Aruncând ancorele, le-au lăsat în mare, dezlegând în același timp frânghiile cârmei. Ridicând trincheta în direcția vântului, s-au îndreptat spre plajă.

\VS{41}Dar, ajungând într-un loc unde se întâlneau două mări, au eșuat nava. Prora a lovit și a rămas nemișcată, dar pupa a început să se rupă din cauza violenței valurilor.


\par }{\PP \VS{42}Sfatul ostașilor era să ucidă pe cei prizonieri, ca să nu mai înoate niciunul dintre ei și să scape.

\VS{43}Dar centurionul, vrând să-l salveze pe Pavel, i-a oprit din planul lor și a poruncit ca cei care știau să înoate să se arunce primii peste bord, ca să se îndrepte spre uscat;

\VS{44}iar ceilalți să îi urmeze, unii pe scânduri și alții pe alte lucruri din corabie. Astfel au scăpat cu toții teferi spre uscat.



\par }\Chap{28}{\PP \VerseOne{1}După ce am scăpat, au
\FTNT{*}{{\FR 28:1 }{\FT{NU citește “noi”}}}aflat că insula se numea Malta.

\VS{2}Băștinașii ne-au arătat o bunătate ieșită din comun, căci au aprins un foc și ne-au primit pe toți, din cauza ploii de acum și din cauza frigului.

\VS{3}Dar când Pavel a strâns un mănunchi de bețe și le-a pus pe foc, o viperă a ieșit din cauza căldurii și s-a prins de mâna lui.

\VS{4}Când băștinașii au văzut creatura atârnând de mâna lui, și-au spus unii altora: “Fără îndoială că acest om este un ucigaș, pe care, deși a scăpat din mare, Justiția nu l-a lăsat să trăiască.”

\VS{5}Cu toate acestea, el a scuturat creatura în foc și nu a pățit nimic.

\VS{6}Ei se așteptau însă ca el să se fi umflat sau să fi căzut mort deodată, dar când au privit mult timp și au văzut că nu i s-a întâmplat nimic rău, s-au răzgândit și au spus că este un zeu.


\par }{\PP \VS{7}În împrejurimile acelui loc se aflau niște pământuri ale căpeteniei insulei, numită Publius, care ne-a primit și ne-a găzduit cu amabilitate timp de trei zile.

\VS{8}Tatăl lui Publius zăcea bolnav de febră și de dizenterie. Pavel a intrat la el, s-a rugat și, punându-și mâinile peste el, l-a vindecat.

\VS{9}După ce s-a făcut acest lucru, au venit și ceilalți care aveau boli în insulă și s-au vindecat.

\VS{10}De asemenea, ne-au onorat cu multe onoruri și, când am plecat, au pus la bord lucrurile de care aveam nevoie.


\par }{\PP \VS{11}După trei luni, am plecat cu o corabie din Alexandria care iernase pe insulă și care avea ca emblemă “Frații gemeni”.

\VS{12}Aterizând la Siracuza, am rămas acolo trei zile.

\VS{13}De acolo am făcut un ocol și am ajuns la Rhegium. După o zi, a apărut un vânt dinspre sud și, a doua zi, am ajuns la Puteoli,

\VS{14}unde am găsit frați și am fost rugați să rămânem cu ei timp de șapte zile. Astfel am ajuns la Roma.

\VS{15}De acolo, frații, când au auzit de noi, au venit în întâmpinarea noastră până la Piața lui Appius și la Cele Trei Taverne. Când i-a văzut, Pavel a mulțumit lui Dumnezeu și a prins curaj.

\VS{16}Când am intrat în Roma, centurionul i-a predat pe prizonieri căpitanului gărzii, dar lui Pavel i s-a permis să rămână singur cu soldatul care îl păzea.


\par }{\PP \VS{17}După trei zile, Pavel a chemat laolaltă pe cei care erau conducătorii iudeilor. După ce s-au adunat, le-a zis: “Eu, fraților, deși nu făcusem nimic împotriva poporului și a obiceiurilor părinților noștri, am fost totuși dat prizonier din Ierusalim în mâinile romanilor,

\VS{18}care, după ce m-au cercetat, au vrut să mă elibereze, pentru că nu era în mine nici o cauză de moarte.

\VS{19}Dar, când iudeii s-au opus, am fost nevoit să apelez la Cezar, nu că aș fi avut ceva de care să acuz națiunea mea.

\VS{20}De aceea am cerut să te văd și să vorbesc cu tine. Căci, din cauza speranței lui Israel, sunt legat cu acest lanț.”


\par }{\PP \VS{21}Ei I-au zis: “Nu am primit scrisori din Iudeea despre tine, și nici unul din frați n-a venit aici să ne spună sau să ne vorbească de rău despre tine.

\VS{22}Dar noi dorim să auzim de la tine ce crezi tu. Căci, în ceea ce privește această sectă, ne este cunoscut faptul că peste tot se vorbește împotriva ei.”


\par }{\PP \VS{23}După ce I-au dat o zi, a venit multă lume la El, la locuința lui. El le dădea explicații, mărturisind despre Împărăția lui Dumnezeu și convingându-i despre Isus, atât din Legea lui Moise, cât și din profeți, de dimineața până seara.

\VS{24}Unii credeau cele spuse, iar alții nu credeau.

\VS{25}Când nu s-au înțeles între ei, au plecat după ce Pavel a rostit un singur mesaj: “Duhul Sfânt a vorbit cu dreptate părinților noștri, prin profetul Isaia,

\VS{26}spunând:

\par }{\Q “Du-te la acest popor și spune-i,

\par }{\Q în auz, veți auzi,

\par }{\QB dar nu va înțelege în nici un fel.

\par }{\Q Văzând, veți vedea,

\par }{\QB dar nu va percepe în nici un fel.


\par }{\Q \VS{27}Căci inima acestui popor a devenit insensibilă.

\par }{\QB Urechile lor sunt surde de auz.

\par }{\QB Și-au închis ochii.

\par }{\Q Ca nu cumva să vadă cu ochii lor,

\par }{\QB aud cu urechile lor,

\par }{\QB să înțeleagă cu inima lor,

\par }{\QB și s-ar întoarce din nou,

\par }{\QB atunci îi voi vindeca.
\FTNT{†}{{\FR 28:27 }{\FT{{\CHAR{Isaia 6:9-10}}}}}


\par }{\PP \VS{28}“Să știți deci că mântuirea lui Dumnezeu este trimisă la neamuri și ele vor asculta.”


\par }{\PP \VS{29}După ce a spus aceste cuvinte, iudeii au plecat, având între ei o mare ceartă.
\FTNT{‡}{{\FR 28:29 }{\FT{NU omite versetul 29.}}}


\par }{\PP \VS{30}Pavel a stat doi ani întregi în casa lui închiriată, și primea pe toți cei ce veneau la el,

\VS{31}propovăduind cu toată îndrăzneala, fără nici o opreliște, Împărăția lui Dumnezeu și învățând cele despre Domnul Isus Hristos.


\par }